% !TEX program = xelatex
\PassOptionsToPackage{unicode=true}{hyperref} % options for packages loaded elsewhere
\PassOptionsToPackage{hyphens}{url}
%
\documentclass[UTF8]{ctexart}
\usepackage{lmodern}
\usepackage{amssymb,amsmath}
\usepackage{ifxetex,ifluatex}
\usepackage{fixltx2e} % provides \textsubscript
\ifnum 0\ifxetex 1\fi\ifluatex 1\fi=0 % if pdftex
  \usepackage[T1]{fontenc}
  \usepackage[utf8]{inputenc}
  \usepackage{textcomp} % provides euro and other symbols
\else % if luatex or xelatex
  \usepackage{unicode-math}
  \defaultfontfeatures{Ligatures=TeX,Scale=MatchLowercase}
\fi
% use upquote if available, for straight quotes in verbatim environments
\IfFileExists{upquote.sty}{\usepackage{upquote}}{}
% use microtype if available
\IfFileExists{microtype.sty}{%
\usepackage[]{microtype}
\UseMicrotypeSet[protrusion]{basicmath} % disable protrusion for tt fonts
}{}
\IfFileExists{parskip.sty}{%
\usepackage{parskip}
}{% else
\setlength{\parindent}{0pt}
\setlength{\parskip}{6pt plus 2pt minus 1pt}
}
\usepackage{hyperref}
\hypersetup{
            pdfborder={0 0 0},
            breaklinks=true}
\urlstyle{same}  % don't use monospace font for urls
\usepackage[margin=1in]{geometry}
\usepackage{longtable,booktabs}
% Fix footnotes in tables (requires footnote package)
\IfFileExists{footnote.sty}{\usepackage{footnote}\makesavenoteenv{longtable}}{}
\usepackage{graphicx,grffile}
\makeatletter
\def\maxwidth{\ifdim\Gin@nat@width>\linewidth\linewidth\else\Gin@nat@width\fi}
\def\maxheight{\ifdim\Gin@nat@height>\textheight\textheight\else\Gin@nat@height\fi}
\makeatother
% Scale images if necessary, so that they will not overflow the page
% margins by default, and it is still possible to overwrite the defaults
% using explicit options in \includegraphics[width, height, ...]{}
\setkeys{Gin}{width=\maxwidth,height=\maxheight,keepaspectratio}
\setlength{\emergencystretch}{3em}  % prevent overfull lines
\providecommand{\tightlist}{%
  \setlength{\itemsep}{0pt}\setlength{\parskip}{0pt}}
\setcounter{secnumdepth}{0}
% Redefines (sub)paragraphs to behave more like sections
\ifx\paragraph\undefined\else
\let\oldparagraph\paragraph
\renewcommand{\paragraph}[1]{\oldparagraph{#1}\mbox{}}
\fi
\ifx\subparagraph\undefined\else
\let\oldsubparagraph\subparagraph
\renewcommand{\subparagraph}[1]{\oldsubparagraph{#1}\mbox{}}
\fi

% set default figure placement to htbp
\makeatletter
\def\fps@figure{htbp}
\makeatother


\date{}

\begin{document}

\hypertarget{sd-mctsux9762ux5411ux8fd0ux7b79ux5b66ux7684ux7ed3ux6784ux5316ux4e16ux754cux6a21ux578b}{%
\section{SD-MCTS：面向运筹学的结构化世界模型}\label{sd-mctsux9762ux5411ux8fd0ux7b79ux5b66ux7684ux7ed3ux6784ux5316ux4e16ux754cux6a21ux578b}}

\textbf{目标期刊}：INFORMS Journal on Computing (IJOC)

\begin{center}\rule{0.5\linewidth}{0.4pt}\end{center}

\hypertarget{ux6458ux8981}{%
\subsection{摘要}\label{ux6458ux8981}}

运筹系统中的序贯决策常受困于延迟效应------当因果反馈滞后于决策周期时，短视界规划难以捕捉动态，而长视界规划在计算上难以承受。这一困境在供应链、库存管理和流程控制等领域普遍存在，却缺乏原则性解决方案。尽管无模型强化学习展现了端到端决策的潜力，其样本需求与可解释性缺失使其难以融入追求结构化建模与可验证性的运筹实践。

我们提出SD-MCTS，一个融合三项设计贡献的混合规划框架：（1）将系统动力学存量-流量语义嵌入为世界模型的硬架构先验，确保守恒约束并实现小样本学习；（2）以学习型价值函数作为终端估值，使浅层搜索即可达成近似无限视界最优；（3）揭示一条实践者决策边界------当特征延迟超过可行规划视界两倍时（\(\tau > 2H\)），投资价值学习优于参数校准，因为视界截断误差以两个数量级主导参数估计误差（\(p < 0.0001\)）。

实验验证指出：在合成延迟崩溃任务中，SD-MCTS以穷举搜索3\%的计算成本达到等效解质量；在啤酒游戏供应链基准上，SD-MCTS显著减少了跨需求场景的手动视界调优需求。结果表明：领域结构可将样本效率提升一个数量级，而终端价值函数为打破延迟-视界权衡提供了可扩展路径。

\textbf{关键词}：系统动力学、蒙特卡洛树搜索、近似动态规划、供应链管理、基于模型的强化学习

\begin{center}\rule{0.5\linewidth}{0.4pt}\end{center}

\hypertarget{ux5f15ux8a00}{%
\subsection{1. 引言}\label{ux5f15ux8a00}}

\hypertarget{ux52a8ux673aux52a8ux6001ux8fd0ux8425ux4e2dux7684ux5ef6ux8fdfux6548ux5e94ux6311ux6218}{%
\subsubsection{1.1
动机：动态运营中的延迟效应挑战}\label{ux52a8ux673aux52a8ux6001ux8fd0ux8425ux4e2dux7684ux5ef6ux8fdfux6548ux5e94ux6311ux6218}}

许多运营系统都呈现出一个典型特征：今天做出的决策，其后果往往在相当长的延迟后才会显现。仓库工人的任务分配不仅影响即时吞吐量，还会导致累积疲劳从而降低未来生产力。机器的运行强度不仅影响当前产出，还会造成磨损累积最终触发故障。服务系统的人员配置不仅影响当前队列长度，还会导致员工职业倦怠从而引发未来人员流失。

这些系统共享三个结构性特征，对传统优化方法构成挑战：

\begin{enumerate}
\def\labelenumi{\arabic{enumi}.}
\item
  \textbf{混合动态}：离散控制决策（分配/不分配、运行/闲置）与连续状态演化（疲劳累积、磨损进展、队列动态）相互作用。
\item
  \textbf{延迟后果}：决策的影响通过累积状态变量传播------系统动力学（SD）称之为``存量''------在行动与其最终结果之间产生时间差。
\item
  \textbf{结构-参数不对称}：领域专家通常能够指定控制方程的\emph{形式}（守恒律、退化函数），但精确\emph{参数}（疲劳率、故障阈值）难以校准。
\end{enumerate}

\hypertarget{ux6edaux52a8ux4f18ux5316ux4e2dux7684ux89c6ux754cux95eeux9898}{%
\paragraph{滚动优化中的视界问题}\label{ux6edaux52a8ux4f18ux5316ux4e2dux7684ux89c6ux754cux95eeux9898}}

模型预测控制（MPC）和滚动时域方法是运筹学中动态优化的主力工具。它们求解有限时域子问题并执行第一个动作，随时间推进重复此过程。这种方法计算上可行且能自然处理约束。

然而，MPC存在一个根本性限制：\textbf{视界截断偏差}。当规划视界 \(H\)
短于系统的特征时间常数
\(\tau\)------即累积效应显现的时间尺度------MPC会系统性地低估那些牺牲短期性能换取长期收益的行动。

考虑一个具体例子：如果疲劳诱发的崩溃发生在第12步，但MPC仅规划5步，它无法``看到''即将到来的故障。结果是短视策略最大化即时回报，却驶向最终灾难。

\hypertarget{ux7eafux9ed1ux76d2ux9884ux6d4bux7684ux4e0dux53efux7ea6ux504fux5deeux4e0eux957fux671fux89c4ux5212ux5931ux6548}{%
\paragraph{纯黑盒预测的不可约偏差与长期规划失效}\label{ux7eafux9ed1ux76d2ux9884ux6d4bux7684ux4e0dux53efux7ea6ux504fux5deeux4e0eux957fux671fux89c4ux5212ux5931ux6548}}

在另一个极端，无模型强化学习（RL）方法（如
PPO、SAC）或端到端的纯黑盒模型基方法（如
Dreamer、MuZero）原则上可以学习优化无限时域目标并提高样本效率。然而，它们在面临具有深度延迟的运筹系统时遭遇了一个致命的理论硬伤：\textbf{长期收益估计的不可约偏差
(Irreducible Bias)}。

当直接依赖缺乏物理结构底座的黑盒神经网络进行多步前置推演时，其状态预测误差会随视界深度
\(H\)
呈指数级放大。这造成了一个结构性的灾难：即使无限增加马尔可夫决策搜索的模拟采样次数，由模型本身预测外溢引发的
\(Q\) 值估计误差 \(\mathcal{O}((L)^H)\) 也永远无法消除（其中 \(L>1\)
为网络局部 Lipschitz
常数）。结果是黑盒模型不仅在训练分布外极易导向失效策略，其在执行安全关键的长周期预测时，更可能幻视出突破系统底层物理约束（如负库存、虚假供给）的错误图景，丧失运筹决策所需的可信度。

\hypertarget{ux8d21ux732e}{%
\subsubsection{1.2 贡献}\label{ux8d21ux732e}}

本文提出\textbf{SD-MCTS}（系统动力学-蒙特卡洛树搜索），一个混合规划框架，弥合了结构感知运筹学与数据驱动机器学习之间的差距。我们的贡献如下：

\textbf{贡献1：面向运筹学的结构化世界模型集成}

我们将系统动力学（SD）的存量-流量建模范式与蒙特卡洛树搜索规划相集成。核心思想是\emph{结构-参数分离}------这一概念在物理信息神经网络（PINNs）、灰箱系统辨识和哈密顿网络中已有研究。我们的贡献在于将这一思想\textbf{具体应用于OR领域的序贯决策问题}：
-
\textbf{结构}（存量-流量方程、守恒律、约束形式）由OR领域知识指定并硬编码。
- \textbf{参数}（速率、阈值、系数）通过梯度下降从数据中学习。

这种集成使样本复杂度从表格方法的
\(O(|\mathcal{S}||\mathcal{A}|/\epsilon^2)\) 降低到
\(O(d/\epsilon^2)\)，其中 \(d\)
是可学习参数的数量------通常是2-5个标量而非数千个状态-动作对。

\textbf{贡献2：用于有效视界扩展的学习终端成本}

我们证明，通过时序差分学习训练的价值网络 \(V_\psi(s)\)
可以作为近似动态规划（ADP）传统中的\emph{学习终端成本}。当用于评估MCTS叶节点时，它将无限时域尾部压缩为单个标量，有效扩展规划视界而无需成比例增加计算成本。

实验表明，搜索深度 \(H=1\)
的SD-MCTS在延迟崩溃任务上达到100\%成功率，而零终端成本的MPC需要
\(H \geq 8\)------\textbf{有效视界扩展至少7步}。

\textbf{贡献3：面向实践者的计算洞见与决策边界}

我们系统分析了规划深度与模型精度之间的性能-预算权衡，并提供可操作的决策边界。我们的核心发现------经理论和实验双重验证------是：

\begin{quote}
\textbf{在延迟效应系统中（\(\tau > H\)），最小化视界截断偏差比最小化参数误差更为关键。}
\end{quote}

具体而言，参数不完美（25\%估计误差）的SD-MCTS以3.0\%的优势超越参数完美的Oracle
MPC（\(p < 0.0001\)，Cohen's \(d = 2.05\)），同时仅使用3.2\%的计算预算。

\textbf{实践者决策规则}：当系统延迟效应的时间尺度 \(\tau\)
超过MPC规划视界 \(H\)
两倍以上时（\(\tau/H > 2\)），应优先投资于价值函数训练而非参数校准。这一规则在三个测试场景中均得到验证（§7）。

\hypertarget{ux8bbaux6587ux7ec4ux7ec7}{%
\subsubsection{1.3 论文组织}\label{ux8bbaux6587ux7ec4ux7ec7}}

本文其余部分组织如下：

\begin{itemize}
\tightlist
\item
  \textbf{§2
  问题类}：正式定义具有领域动态的结构化MDP类，包括``结构已知、参数未知''假设。
\item
  \textbf{§3
  相关文献}：将SD-MCTS定位于MPC、ADP、基于模型的RL和物理信息学习的研究版图中。
\item
  \textbf{§4
  方法}：介绍SD-MCTS框架，包括结构化世界模型、带学习终端值的MCTS和计算复杂度分析。
\item
  \textbf{§5
  理论分析}：提供样本复杂度、有效视界扩展和视界-精度权衡的理论界限。
\item
  \textbf{§6 计算研究}：描述实验设计、测试平台和基准方法。
\item
  \textbf{§7 结果}：展示主要性能比较、机制研究和可扩展性分析。
\item
  \textbf{§8 可复现性}：描述软件仓库和复现说明。
\item
  \textbf{§9 结论}：总结洞见并讨论未来方向。
\end{itemize}

\hypertarget{ux53efux590dux73b0ux6027ux58f0ux660e}{%
\subsubsection{1.4
可复现性声明}\label{ux53efux590dux73b0ux6027ux58f0ux660e}}

根据INFORMS Journal on Computing软件和数据政策，我们提供完整的复现包：

\begin{verbatim}
https://github.com/[anonymized]/sd-mcts-ijoc
\end{verbatim}

仓库包含： - 所有算法和环境的源代码（\texttt{src/}） -
复现所有表格和图形的脚本（\texttt{scripts/}） -
原始实验结果（\texttt{results/}） -
环境规格（\texttt{requirements.txt}）

所有实验均可通过提供的脚本完全复现；详见§8。

\begin{center}\rule{0.5\linewidth}{0.4pt}\end{center}

\hypertarget{ux95eeux9898ux7c7bux5177ux6709ux9886ux57dfux52a8ux6001ux7684ux7ed3ux6784ux5316mdp}{%
\subsection{2.
问题类：具有领域动态的结构化MDP}\label{ux95eeux9898ux7c7bux5177ux6709ux9886ux57dfux52a8ux6001ux7684ux7ed3ux6784ux5316mdp}}

我们考虑一类序贯决策问题，其中系统动态遵循已知的结构形式但涉及未知参数。这种设定在运筹学中普遍存在：领域专家可以基于第一性原理指定控制方程（队列动态、库存平衡、退化规律），但精确系数需要从数据中估计。

\hypertarget{ux9a6cux5c14ux53efux592bux51b3ux7b56ux8fc7ux7a0bux5efaux6a21}{%
\subsubsection{2.1
马尔可夫决策过程建模}\label{ux9a6cux5c14ux53efux592bux51b3ux7b56ux8fc7ux7a0bux5efaux6a21}}

我们将问题建模为折扣无限时域马尔可夫决策过程（MDP）：

\[\mathcal{M} = (\mathcal{S}, \mathcal{A}, P, R, \gamma, \rho_0)\]

\begin{table}[h]
\centering
\begin{tabular}{p{3cm}p{3cm}p{8cm}}
\toprule
\textbf{符号} & \textbf{定义} & \textbf{描述} \\
\midrule
$\mathcal{S} \subseteq \mathbb{R}^n$ & 状态空间 & 连续，$n$维 \\
$\mathcal{A}$ & 动作空间 & 离散，$\lvert\mathcal{A}\rvert = K$ \\
$P: \mathcal{S} \times \mathcal{A} \to \Delta(\mathcal{S})$ & 转移核 & $P(s'\lvert s,a)$ 为下一状态分布 \\
$R: \mathcal{S} \times \mathcal{A} \to \mathbb{R}$ & 奖励函数 & 有界，$\lvert R(s,a)\rvert \leq R_{\max}$ \\
$\gamma \in (0,1)$ & 折扣因子 & 控制未来奖励权重 \\
$\rho_0$ & 初始分布 & $s_0 \sim \rho_0$ \\
\bottomrule
\end{tabular}
\end{table}

\textbf{目标}：找到最优策略
\(\pi^*: \mathcal{S} \to \Delta(\mathcal{A})\)，最大化期望累积折扣回报：

\[\pi^* = \arg\max_\pi \mathbb{E}_\pi \left[ \sum_{t=0}^{\infty} \gamma^t R(s_t, a_t) \mid s_0 \sim \rho_0 \right]\]

\hypertarget{ux7ed3ux6784-ux53c2ux6570ux5206ux79bbux5047ux8bbe}{%
\subsubsection{2.2
结构-参数分离假设}\label{ux7ed3ux6784-ux53c2ux6570ux5206ux79bbux5047ux8bbe}}

我们问题类的定义特征是转移核具有已知的参数化形式：

\begin{quote}
\textbf{假设1（结构知识）}：转移核具有由
\(\theta \in \Theta \subseteq \mathbb{R}^d\) 参数化的已知函数形式：
\[P_\theta(s'|s,a) = \mathcal{N}(f(s,a;\theta), \Sigma)\] 其中结构函数
\(f: \mathcal{S} \times \mathcal{A} \times \Theta \to \mathcal{S}\)
是已知的（由领域专业知识指定），只有参数向量 \(\theta\) 是未知的。
\end{quote}

\begin{quote}
\textbf{假设2（可实现性）}：真实参数
\(\theta^* \in \Theta\)，即模型类包含真实系统。
\end{quote}

\begin{quote}
\textbf{假设3（可辨识性）}：映射 \(\theta \mapsto P_\theta\)
是单射------不同参数产生可区分的动态。
\end{quote}

\begin{quote}
\textbf{假设4（持续激励）}：数据分布 \(\mathcal{D}\) 满足：
\[\lambda_{\min}\left(\mathbb{E}_{(s,a)\sim\mathcal{D}}\left[\nabla_\theta f \cdot \nabla_\theta f^\top\right]\right) \geq \lambda_0 > 0\]
\end{quote}

\textbf{运筹学应用中的合理性说明}：

\begin{table}[h]
\centering
\begin{tabular}{p{3cm}p{11cm}}
\toprule
\textbf{假设} & \textbf{运筹学背景下的合理性} \\
\midrule
A1（结构） & 存量-流量方程、排队公式和退化规律源自第一性原理；\emph{形式}由领域专业知识确定 \\
A2（可实现性） & 参数空间 \(\Theta\) 由物理约束定义（如 \(\alpha > 0\)，\(\tau > 0\)），确保真实系统位于其中 \\
A3（可辨识性） & SD方程对参数连续可微；不同的 \((\alpha, \tau)\) 产生可区分的轨迹 \\
A4（激励） & MCTS探索（PUCT）系统性覆盖状态-动作空间；结合可观测存量，Fisher信息满秩 \\
\bottomrule
\end{tabular}
\end{table}

\hypertarget{ux5b58ux91cf-ux6d41ux91cfux52a8ux6001ux8fd0ux7b79ux5b66ux7cfbux7edfux7684ux6a21ux677f}{%
\subsubsection{2.3
存量-流量动态：运筹学系统的模板}\label{ux5b58ux91cf-ux6d41ux91cfux52a8ux6001ux8fd0ux7b79ux5b66ux7cfbux7edfux7684ux6a21ux677f}}

系统动力学（Forrester
1961）为在运营中表达结构化动态提供了自然语言。基本构建块是\textbf{存量-流量}关系：

\[\frac{dS(t)}{dt} = \text{流入}(t) - \text{流出}(t)\]

这种守恒结构广泛适用于各运筹学领域：

\begin{table}[h]
\centering
\begin{tabular}{p{3cm}p{5cm}p{6cm}}
\toprule
\textbf{领域} & \textbf{存量变量} & \textbf{流量变量} \\
\midrule
库存 & 现有库存、缺货订单 & 补货、需求、损耗 \\
劳动力 & 可用产能、疲劳水平 & 任务完成、休息、人员流失 \\
排队 & 队列长度、服务器利用率 & 到达、离开、放弃 \\
维护 & 设备健康度、磨损水平 & 使用退化、维修恢复 \\
\bottomrule
\end{tabular}
\end{table}

\textbf{离散时间建模}：为计算可行性，我们以时间步 \(\Delta t\) 离散化：

\[S_{t+1} = S_t + (\text{流入}_t - \text{流出}_t) \cdot \Delta t\]

\textbf{示例：带疲劳的任务分配}（我们的主要测试平台）：

\[q^{(t+1)} = \max\left(0, q^{(t)} + A^{(t)} - \textstyle\sum_i x_i^{(t)}\right)\]
\[b_i^{(t+1)} = \max\left(0, b_i^{(t)} + x_i^{(t)} - o_i^{(t)}\right)\]
\[f^{(t+1)} = (1-\rho) f^{(t)} + \kappa \textstyle\sum_i o_i^{(t)}\]

其中： - \(q^{(t)}\)：队列长度（未分配任务的存量） - \(b_i^{(t)}\)：工人 \(i\)
的积压（已分配但未完成任务的存量） -
\(f^{(t)}\)：累积疲劳（累积工作负荷效应的存量） -
\(A^{(t)} \sim \text{Poisson}(\lambda \Delta t)\)：随机到达（流入） -
\(x_i^{(t)}\)：分配决策（受控流量） -
\(o_i^{(t)} = \eta(f^{(t)}) \cdot \mu_i\)：有效服务率（流出），受疲劳调节

\textbf{可学习参数}：\(\theta = (\alpha, \tau, \lambda, \mu_1, \ldots, \mu_n)\)

\begin{itemize}
\tightlist
\item
  \(\alpha\)：疲劳敏感度（退化率）
\item
  \(\tau\)：疲劳阈值（退化前的容忍度）
\item
  \(\lambda\)：到达率
\item
  \(\mu_i\)：工人 \(i\) 的基础服务率
\end{itemize}

\hypertarget{ux5904ux7406ux968fux673aux6027ux5e73ux5747ux573aux8fd1ux4f3c}{%
\subsubsection{2.4
处理随机性：平均场近似}\label{ux5904ux7406ux968fux673aux6027ux5e73ux5747ux573aux8fd1ux4f3c}}

存量-流量方程涉及随机元素（如泊松到达
\(A_t\)）。为实现可微分学习，我们采用\textbf{平均场近似}：

\[A_t \approx \mathbb{E}[A_t] = \lambda \Delta t\]

这用期望值替代离散随机事件，在保持结构形式的同时实现基于梯度的参数学习。

\begin{quote}
\textbf{备注（与流体近似的联系）}：我们的平均场方法与排队论中的\emph{流体近似}（Mandelbaum
1998，Whitt
2002）有哲学上的共同根源，后者通过大数定律标度将离散随机过程（如
\(M/M/1\)
队列）近似为确定性微分方程。系统动力学对具有累积和延迟结构的反馈系统采用类似的建模精神。我们强调这是\emph{概念性}联系而非形式等价------我们不假设标度机制或引用收敛结果，该近似纯粹用作基于梯度学习的可微代理。
\end{quote}

\textbf{权衡}：

\begin{table}[h]
\centering
\begin{tabular}{p{3cm}p{5cm}p{6cm}}
\toprule
\textbf{方面} & \textbf{平均场} & \textbf{真实随机} \\
\midrule
可微性 & ✓ 端到端梯度 & ✗ 需要REINFORCE/重参数化 \\
训练稳定性 & ✓ 较低方差 & ✗ 较高方差 \\
决策质量 & 轻微乐观偏差 & 精确 \\
计算成本 & 较低 & 较高 \\
\bottomrule
\end{tabular}
\end{table}

\begin{quote}
\textbf{备注}：对于规划（MCTS），我们可以选择性地在树回滚期间通过采样
\(A_t \sim \text{Poisson}(\lambda \Delta t)\)
恢复随机性，同时使用平均场进行基于梯度的模型学习。我们的实验（§7）表明这种混合方法与完全随机模拟性能相当，且计算成本更低。
\end{quote}

\begin{quote}
\textbf{备注（与理论结果的关系）}：§5中的理论界限（引理1、定理1-2）假设可访问真实转移动态。在平均场近似下，存在额外偏差项
\(O(\text{Var}[A_t]/\mathbb{E}[A_t])\)，对于高速率泊松过程通常较小。实验上，这种偏差不影响定性结论：视界扩展机制（定理1）和视界-精度权衡（定理2）按预测运作。
\end{quote}

\hypertarget{ux4ef7ux503cux51fdux6570ux4e0eux8d1dux5c14ux66fcux65b9ux7a0b}{%
\subsubsection{2.5
价值函数与贝尔曼方程}\label{ux4ef7ux503cux51fdux6570ux4e0eux8d1dux5c14ux66fcux65b9ux7a0b}}

为完整性，我们定义本方法核心的价值函数：

\textbf{状态价值函数}：
\[V^\pi(s) = \mathbb{E}_\pi\left[\sum_{t=0}^{\infty} \gamma^t R(s_t, a_t) \mid s_0 = s\right]\]

\textbf{动作价值函数}：
\[Q^\pi(s,a) = R(s,a) + \gamma \mathbb{E}_{s' \sim P(\cdot|s,a)}[V^\pi(s')]\]

\textbf{贝尔曼最优方程}：
\[V^*(s) = \max_{a \in \mathcal{A}} \left\{R(s,a) + \gamma \mathbb{E}_{s'}[V^*(s')]\right\}\]

\textbf{截断价值（MPC视角）}：
\[V_H^\pi(s) = \mathbb{E}_\pi\left[\sum_{t=0}^{H-1} \gamma^t R(s_t, a_t) + \gamma^H \cdot V_{\text{terminal}}(s_H)\right]\]

\(V_{\text{terminal}}\) 的选择至关重要： -
\textbf{MPC}：\(V_{\text{terminal}} = 0\)（零终端成本）→ 截断偏差
\(O(\gamma^H R_{\max}/(1-\gamma))\) -
\textbf{SD-MCTS}：\(V_{\text{terminal}} = V_\psi(s_H)\)（学习价值）→
截断偏差 \(O(\gamma^H \epsilon_V)\)

\begin{center}\rule{0.5\linewidth}{0.4pt}\end{center}

\hypertarget{ux76f8ux5173ux6587ux732e}{%
\subsection{3. 相关文献}\label{ux76f8ux5173ux6587ux732e}}

SD-MCTS汲取了四个研究流派：滚动时域优化、近似动态规划、基于模型的强化学习和物理信息机器学习。我们在每个传统中定位本工作。

\hypertarget{ux6edaux52a8ux65f6ux57dfux4f18ux5316ux4e0eux6a21ux578bux9884ux6d4bux63a7ux5236}{%
\subsubsection{3.1
滚动时域优化与模型预测控制}\label{ux6edaux52a8ux65f6ux57dfux4f18ux5316ux4e0eux6a21ux578bux9884ux6d4bux63a7ux5236}}

模型预测控制（MPC）是运筹学和控制领域动态优化的主导范式（Camacho and
Bordons 2007, Rawlings et
al.~2017）。在每个决策时刻，MPC求解一个有限时域优化问题并执行第一个动作，然后用更新的状态信息重新求解。

\textbf{MPC中的终端成本}：缓解视界截断的经典方法是附加一个\emph{终端成本}
\(V_{\text{terminal}}(s_H)\)，近似视界之外的价值函数（Mayne et
al.~2000）。常见选择包括： -
零终端成本：\(V_{\text{terminal}} = 0\)（最简单，但引入截断偏差） -
二次近似：\(V_{\text{terminal}} = s_H^\top P s_H\)（需要求解Riccati方程）
- 控制Lyapunov函数：保证稳定性但需要显式构造

\textbf{我们的贡献}：SD-MCTS可以被视为具有\emph{学习}终端成本的MPC。我们不假设参数形式（二次）或求解解析表达式，而是通过时序差分学习训练神经网络
\(V_\psi\) 来近似
\(V^*\)。这将经典终端成本思想扩展到解析解不可用的场景。

\hypertarget{ux8fd1ux4f3cux52a8ux6001ux89c4ux5212ux4e0eux56deux6edaux65b9ux6cd5}{%
\subsubsection{3.2
近似动态规划与回滚方法}\label{ux8fd1ux4f3cux52a8ux6001ux89c4ux5212ux4e0eux56deux6edaux65b9ux6cd5}}

近似动态规划（ADP）通过近似价值函数来解决动态规划中的维数灾难（Bertsekas
2019, Powell 2011）。关键方法论包括：

\textbf{价值函数近似}：使用基函数 \(\phi(s)\) 表示
\(V(s) \approx \phi(s)^\top w\)，并通过模拟学习权重 \(w\)（Tsitsiklis
and Van Roy 1997）。

\textbf{回滚与策略改进}：从基础策略 \(\pi_0\) 出发，回滚通过基于
\(\pi_0\) 评估的单步前瞻计算改进策略（Bertsekas and Tsitsiklis
1996）。这等价于策略迭代中的策略改进。

\textbf{与SD-MCTS的关系}：我们的框架是\emph{带学习终端评估的多步回滚}的一种形式：
- 经典回滚：1步前瞻 + 基础策略回滚至视界 - SD-MCTS：\(H\)步树搜索 +
\(V_\psi\) 终端评估

价值网络 \(V_\psi\)
在ADP术语中扮演\emph{终端成本函数}的角色------它将无限时域尾部压缩为单个标量，使有限深度搜索能够考虑长期后果。

\begin{quote}
\textbf{澄清}：我们不声称在使用学习价值函数进行规划的\emph{思想}上有新颖性------这在ADP中已有充分建立。我们的贡献在于(i)将其与编码领域物理的\emph{结构化}世界模型结合，以及(ii)展示学习终端成本在延迟效应系统中补偿模型不完美的\emph{机制}。
\end{quote}

\hypertarget{ux57faux4e8eux6a21ux578bux7684ux5f3aux5316ux5b66ux4e60ux4e0eux4e16ux754cux6a21ux578b}{%
\subsubsection{3.3
基于模型的强化学习与世界模型}\label{ux57faux4e8eux6a21ux578bux7684ux5f3aux5316ux5b66ux4e60ux4e0eux4e16ux754cux6a21ux578b}}

基于模型的RL学习环境动态并将其用于规划或策略优化（Moerland et
al.~2023）。近期进展包括：

\textbf{隐空间世界模型}：Dreamer（Hafner et
al.~2020）和MuZero（Schrittwieser et
al.~2020）在隐空间中学习动态，在游戏和连续控制上取得强劲性能。然而，它们\emph{从零开始}学习动态，不融入结构知识。

\textbf{使用学习模型进行规划}：AlphaZero（Silver et
al.~2017）将MCTS与神经网络策略和价值函数结合，使用已知模拟器进行状态转移。MuZero将此扩展到学习的动态。

\textbf{与SD-MCTS的比较}：

\begin{table}[h]
\centering
\begin{tabular}{p{3cm}p{5cm}p{6cm}}
\toprule
\textbf{方面} & \textbf{MuZero/Dreamer} & \textbf{SD-MCTS} \\
\midrule
动态 & 端到端学习（黑盒） & 结构已知，参数学习 \\
状态空间 & 隐空间（不可解释） & 可观测（物理意义明确） \\
样本效率 & \(\sim 10^6\) 步 & \(\sim 10^4\) 步 \\
约束处理 & 隐式（须学习） & 显式（硬编码） \\
泛化 & 分布偏移下表现差 & 鲁棒（结构保持） \\
\bottomrule
\end{tabular}
\end{table}

\textbf{关键区别}：MuZero学习动态\emph{是什么}；SD-MCTS学习已知动态的\emph{参数}。这种结构-参数分离是我们样本效率优势的来源。

\hypertarget{ux7269ux7406ux4fe1ux606fux4e0eux7ed3ux6784ux4fddux6301ux5b66ux4e60}{%
\subsubsection{3.4
物理信息与结构保持学习}\label{ux7269ux7406ux4fe1ux606fux4e0eux7ed3ux6784ux4fddux6301ux5b66ux4e60}}

越来越多的工作将物理结构融入机器学习：

\textbf{物理信息神经网络（PINNs）}：Raissi et al.
(2019)将微分方程作为软约束嵌入神经网络训练。网络学习在拟合数据的同时满足控制PDEs。

\textbf{神经ODEs}：Chen et al.
(2018)用神经网络参数化连续时间动态，并使用ODE求解器进行前向模拟。这实现了微分方程模型的端到端学习。

\textbf{哈密顿/拉格朗日网络}：Greydanus et al. (2019)和Cranmer et al.
(2020)硬编码能量守恒结构，仅学习哈密顿量或拉格朗日量函数。

\textbf{运筹学中的系统动力学}：SD在运筹学中有悠久的策略分析历史（Forrester
1961, Sterman 2000），但通常使用手动校准参数而非数据驱动学习。

\textbf{为什么不用时序预测？}：结构化世界模型的替代方案是统计预测（ARIMA、指数平滑、Prophet）。然而，时序方法预测\emph{结果}（如未来队列长度）而非\emph{因果动态}（动作如何影响结果）。它们无法回答规划所需的反事实查询：``如果我们分配更少任务会怎样？''SD模型嵌入因果结构，能够模拟假设策略------这是树搜索规划的基本要求。

\textbf{与经典OR方法的联系}：MCTS树搜索与组合优化中的分支定界共享概念DNA：两者都探索决策树，使用界限剪枝无望分支。价值网络
\(V_\psi\)
扮演类似整数规划中对偶界的角色------它提供价值函数估计来指导搜索优先级。类似地，结构化世界模型与随机规划中的场景树相关：两者都通过分支结构表示不确定未来，尽管MCTS实时采样场景而非预先指定。

\textbf{我们的定位}：SD-MCTS在精神上最接近哈密顿网络------我们硬编码\emph{结构}（存量-流量方程）并学习\emph{参数}。然而，我们关注序贯决策而非系统辨识，并将学习模型与树搜索规划集成。

\hypertarget{ux603bux7ed3ux5b9aux4f4dsd-mcts}{%
\subsubsection{3.5
总结：定位SD-MCTS}\label{ux603bux7ed3ux5b9aux4f4dsd-mcts}}

\begin{table}[h]
\centering
\begin{tabular}{p{4cm}p{5cm}p{5cm}}
\toprule
\textbf{研究流派} & \textbf{我们继承的} & \textbf{我们添加的} \\
\midrule
\textbf{MPC} & 滚动优化，终端成本概念 & 通过神经网络\emph{学习}终端成本 \\
\textbf{ADP} & 价值近似，回滚 & \emph{结构化}世界模型 + MCTS \\
\textbf{基于模型的RL} & 使用学习模型规划，MCTS & \emph{物理信息}结构，样本效率 \\
\textbf{物理信息ML} & 结构保持，可解释性 & \emph{决策制定}焦点，树搜索集成 \\
\bottomrule
\end{tabular}
\end{table}

SD-MCTS占据独特位置：它结合了OR模型的结构可解释性与现代RL的学习能力，以及树搜索的规划能力与ADP终端成本的视界扩展。

\hypertarget{ux4ec0ux4e48ux662fux771fux6b63ux7684ux65b0ux8d21ux732e}{%
\paragraph{3.5.1
什么是真正的新贡献？}\label{ux4ec0ux4e48ux662fux771fux6b63ux7684ux65b0ux8d21ux732e}}

为了回应SD-MCTS是否仅是已知技术简单组合的合理质疑，我们澄清单个组件\emph{无法}实现的内容：

\begin{enumerate}
\def\labelenumi{\arabic{enumi}.}
\item
  \textbf{单独的ADP无法利用物理结构。}
  经典ADP使用忽略领域物理的通用基函数（多项式、瓦片编码）。SD-MCTS硬编码守恒律，确保模型预测永不违反物理约束（如非负库存）。这消除了价值函数方法必须学习避免的一整类错误。
\item
  \textbf{单独的MPC无法在不导致计算爆炸的情况下扩展有效视界。}
  将MPC视界从 \(H=5\) 增加到 \(H=15\) 会使穷举搜索的计算成本增加
  \(|A|^{10}\) 倍。SD-MCTS通过 \(O(1)\)
  评估开销的学习终端成本实现等效视界扩展。
\item
  \textbf{单独的黑盒世界模型无法实现样本效率。} MuZero需要 \(\sim 10^6\)
  样本来学习Atari动态。SD-MCTS仅需 \(\sim 10^4\)
  样本，因为它只在已知结构内学习 \(d \ll |\mathcal{S}|\) 个参数。
\item
  \textbf{组合产生涌现能力。}
  任何单独组件都无法发现``主动休息''策略（§7.1.2）：没有结构的ADP无法建模疲劳动态；没有终端价值的MPC无法看到崩溃之后。协同效应就是贡献。
\end{enumerate}

\begin{center}\rule{0.5\linewidth}{0.4pt}\end{center}

\hypertarget{ux65b9ux6cd5sd-mctsux6846ux67b6}{%
\subsection{4.
方法：SD-MCTS框架}\label{ux65b9ux6cd5sd-mctsux6846ux67b6}}

我们介绍SD-MCTS，一个将结构化领域动态与蒙特卡洛树搜索集成的混合规划框架。关键洞见是利用已知物理结构（存量-流量方程）作为归纳偏置，同时仅学习少量未知参数，并使用学习价值网络作为近似动态规划（ADP）意义上的近似终端成本。

\hypertarget{ux6846ux67b6ux6982ux8ff0}{%
\subsubsection{4.1 框架概述}\label{ux6846ux67b6ux6982ux8ff0}}

SD-MCTS由四个紧密集成的组件组成，在迭代循环中运行（见算法1）：

\begin{enumerate}
\def\labelenumi{\arabic{enumi}.}
\tightlist
\item
  \textbf{数据收集}：在环境中执行当前MCTS策略，收集状态-动作-奖励轨迹。
\item
  \textbf{世界模型学习}：通过状态预测损失上的梯度下降更新结构化动态参数
  \(\theta\)。
\item
  \textbf{价值网络学习}：使用时序差分（TD）学习训练 \(V_\psi(s)\) 以近似
  \(V^*(s)\)。
\item
  \textbf{策略蒸馏}（可选）：将MCTS访问计数蒸馏到策略网络
  \(\pi_\phi(a|s)\) 以加速动作选择。
\end{enumerate}

\textbf{算法1：SD-MCTS训练循环}

\begin{quote}
\ttfamily
输入：环境 env，初始参数 $\theta_0$, $\psi_0$ \\
输出：训练好的世界模型 $M_\theta$，价值网络 $V_\psi$ \\
~\\
for episode = 1 to N do \\
\hspace*{4em}// 阶段1：数据收集 \\
\hspace*{4em}trajectory $\leftarrow$ CollectEpisode(env, MCTS\_policy($M_\theta$, $V_\psi$)) \\
\hspace*{4em}buffer.add(trajectory) \\
~\\
\hspace*{4em}// 阶段2：世界模型更新（可微分系统辨识） \\
\hspace*{4em}从buffer采样批次 $\{(s_t, a_t, s_{t+1})\}$ \\
\hspace*{4em}L\_model $\leftarrow ||s_{t+1} - M_\theta(s_t, a_t)||^2$ \\
\hspace*{4em}$\theta \leftarrow \theta - \eta_\theta \nabla_\theta \text{L\_model}$ \\
~\\
\hspace*{4em}// 阶段3：价值网络更新（TD学习） \\
\hspace*{4em}L\_value $\leftarrow ||V_\psi(s_t) - (r_t + \gamma \cdot V_\psi(s_{t+1}))||^2$ \\
\hspace*{4em}$\psi \leftarrow \psi - \eta_\psi \nabla_\psi \text{L\_value}$ \\
~\\
\hspace*{4em}// 阶段4：策略蒸馏（可选） \\
\hspace*{4em}$\pi_{\text{target}} \leftarrow \text{normalize}(\text{MCTS\_visit\_counts})$ \\
\hspace*{4em}L\_policy $\leftarrow \text{CrossEntropy}(\pi_\phi(\cdot|s_t), \pi_{\text{target}})$ \\
\hspace*{4em}$\phi \leftarrow \phi - \eta_\phi \nabla_\phi \text{L\_policy}$ \\
end for
\end{quote}

\begin{quote}
\textbf{备注（与AlphaZero/ADP的关系）}：此循环遵循AlphaZero范式（Silver
et
al.~2017），但有两个关键区别：(i)世界模型是\emph{结构化}的而非端到端学习的，(ii)策略网络是可选的------当
\(\pi_\phi\) 不可用时MCTS可以使用均匀先验。在OR术语中，\(V_\psi\)
作为ADP传统中的\emph{学习终端成本函数}（Bertsekas
2019），提供无限时域价值函数的近似。
\end{quote}

\begin{quote}
【图F1占位】框架图示：环境 ↔ 数据收集 → 世界模型学习 → 价值学习 →
MCTS规划 → 动作
\end{quote}

\begin{center}\rule{0.5\linewidth}{0.4pt}\end{center}

\hypertarget{ux7ed3ux6784ux5316ux4e16ux754cux6a21ux578bux4f5cux4e3aux5f52ux7eb3ux504fux7f6eux7684ux7269ux7406ux65b9ux7a0b}{%
\subsubsection{4.2
结构化世界模型：作为归纳偏置的物理方程}\label{ux7ed3ux6784ux5316ux4e16ux754cux6a21ux578bux4f5cux4e3aux5f52ux7eb3ux504fux7f6eux7684ux7269ux7406ux65b9ux7a0b}}

我们不端到端学习任意动态函数
\(s' = f_\theta(s,a)\)，而是使用已知的系统动力学（SD）结构参数化转移模型。以任务分配系统为例：

\textbf{状态表示}：\(s = (q, b_1, \ldots, b_n, f) \in \mathbb{R}^{n+2}\)
- \(q\)：队列长度（等待分配的到达任务） - \(b_i\)：工人 \(i\)
的积压（已分配但未完成的任务） - \(f\)：累积疲劳（捕获延迟退化效应）

\textbf{动作空间}：\(a \in \{0, 1, \ldots, n\}\)（离散分配决策）

\textbf{结构化转移方程}（存量-流量守恒）： \[
q_{t+1} = \text{clamp}\left(q_t + A_t - \sum_i x_{i,t}, \ 0, \ q_{\max}\right)
\] \[
b_i^{(t+1)} = \text{clamp}\left(b_i^{(t)} + x_i^{(t)} - o_i^{(t)}, \ 0, \ b_{\max}\right)
\] \[
f_{t+1} = (1-\rho) f_t + \kappa \sum_i s_{i,t}
\]

其中 \(A_t\) 是随机到达，\(x_{i,t}\)
是分配，\(s_{i,t} = \eta(f_t) \cdot \mu_i\) 是受疲劳依赖效率 \(\eta(f)\)
调节的有效服务率。

\textbf{效率退化函数}（可学习参数 \(\alpha, \tau\)）： \[
\eta(f; \alpha, \tau) = \frac{1}{1 + \alpha \cdot \max(0, f - \tau)^2}
\]

\hypertarget{ux5df2ux77e5ux4e0eux53efux5b66ux4e60ux7ec4ux4ef6}{%
\paragraph{4.2.1
已知与可学习组件}\label{ux5df2ux77e5ux4e0eux53efux5b66ux4e60ux7ec4ux4ef6}}

\begin{table}[h]
\centering
\begin{tabular}{p{3cm}p{2cm}p{4cm}p{5cm}}
\toprule
\textbf{组件} & \textbf{状态} & \textbf{表示} & \textbf{OR解释} \\
\midrule
存量-流量结构 & \textbf{已知} & 硬编码方程 & 守恒律 \\
非负约束 & \textbf{已知} & \(\text{clamp}(\cdot, 0, \max)\) & 库存 \(\geq 0\)，容量界限 \\
到达率 \(\lambda\) & \textbf{可学习} & \(\lambda = \text{softplus}(\phi_\lambda)\) & 需求率 \\
退化率 \(\alpha\) & \textbf{可学习} & \(\alpha = \text{softplus}(\phi_\alpha)\) & 疲劳敏感度 \\
阈值 \(\tau\) & \textbf{可学习} & \(\tau = \text{softplus}(\phi_\tau)\) & 疲劳容忍度 \\
服务率 \(\mu_i\) & \textbf{可学习} & \(\mu_i = 2 \cdot \text{sigmoid}(\phi_\mu)\) & 工人生产力 \\
\bottomrule
\end{tabular}
\end{table}

\begin{quote}
\textbf{关键设计选择}：约束通过clamp操作和有界激活函数\emph{按构造}强制执行，而非通过惩罚项。这保证了所有模型预测的物理可行性（如非负队列）。
\end{quote}

\hypertarget{ux53efux5faeux5206ux7cfbux7edfux8fa8ux8bc6}{%
\paragraph{4.2.2
可微分系统辨识}\label{ux53efux5faeux5206ux7cfbux7edfux8fa8ux8bc6}}

结构化方程形成可微分计算图。参数更新从预测损失通过已知物理结构流动：

\[
\mathcal{L}_{\text{model}} = \|s_{t+1} - \hat{s}_{t+1}\|^2
\]

\[
\theta \leftarrow \theta - \eta \nabla_\theta \mathcal{L}_{\text{model}}
\]

这实现了``通过反向传播的系统辨识''------结构固定；仅调整标量参数。

\hypertarget{ux7ea6ux675fux5904ux7406ux7ed3ux6784ux4fddux8bc1-vs.ux5b66ux4e60ux504fux597d}{%
\paragraph{4.2.3 约束处理：结构保证
vs.~学习偏好}\label{ux7ea6ux675fux5904ux7406ux7ed3ux6784ux4fddux8bc1-vs.ux5b66ux4e60ux504fux597d}}

结构化世界模型相对于黑盒替代方案的一个关键优势是\textbf{硬约束}（由转移模型强制）和\textbf{软偏好}（编码在目标函数中）的分离：

\begin{table}[h]
\centering
\begin{tabular}{p{4cm}p{5cm}p{5cm}}
\toprule
\textbf{约束类型} & \textbf{处理机制} & \textbf{保证级别} \\
\midrule
\textbf{守恒律}（质量平衡、流量连续性） & 硬编码在存量-流量方程中 & \textbf{数学上保证（若已编码）} \\
\textbf{容量限制}（队列上限、库存边界） & 转移函数中的截断操作 & \textbf{构造性强制（若已编码）} \\
\textbf{非负性}（库存 ≥ 0、疲劳 ≥ 0） & 状态更新中的 \(\max(0, \cdot)\) & \textbf{构造性强制（若已编码）} \\
\textbf{时序排序}（\(t\) 时刻状态仅影响 \(t+1\)） & 离散时间更新方案 & \textbf{模型设计保证} \\
\textbf{目标权衡}（服务水平 vs.~成本） & 奖励/惩罚项；由 \(V_\psi\) 近似 & \textbf{非可行性约束} \\
\textbf{多资源可行性}（兼容性、共享容量） & MCTS中的显式动作掩码或可行性检查 & \textbf{仅当显式编码时} \\
\bottomrule
\end{tabular}
\end{table}

\textbf{对样本效率的影响}：与必须通过试错发现物理定律的无模型方法不同（例如，通过奖励惩罚学习库存不能为负），SD-MCTS
\emph{先验地}编码这些约束。学习负担减少为在已知可行结构内为已编码约束估计标量参数，这解释了引理1中的样本复杂度优势。

\begin{quote}
\textbf{备注（预测中的约束违反）}：因为编码的边界是构造性强制的（截断/激活函数），世界模型在其预测中尊重那些\emph{特定}约束。未在转移函数或动作可行性检查中编码的约束仍可能被违反，且截断可能引入建模偏差（例如，队列溢出时的隐式需求损失）。相比之下，黑盒模型（MLPs、LSTMs）甚至可能违反基本边界，需要事后校正或训练时的惩罚项。
\end{quote}

\begin{center}\rule{0.5\linewidth}{0.4pt}\end{center}

\hypertarget{ux5e26ux5b66ux4e60ux7ec8ux7aefux503cux7684mcts}{%
\subsubsection{4.3
带学习终端值的MCTS}\label{ux5e26ux5b66ux4e60ux7ec8ux7aefux503cux7684mcts}}

在决策时，我们使用训练好的世界模型进行状态模拟，使用价值网络进行叶节点评估来运行MCTS。

\textbf{算法2：带学习终端值的MCTS规划}

\begin{verbatim}
function MCTS(s_root, M_θ, V_ψ, N_sim):
    root ← Node(s_root)

    for sim = 1 to N_sim do
        // 选择：使用PUCT遍历树
        node ← root; path ← [root]
        while node.is_expanded() do
            a* ← argmax_a PUCT(node, a)
            node ← node.children[a*]
            path.append(node)
        end while

        // 扩展：使用世界模型预测下一状态
        for each action a ∈ A do
            s' ← M_θ(node.state, a)
            r ← R(node.state, a, s')
            child ← Node(s', prior=π(a|s), reward=r)
            node.children[a] ← child
        end for

        // 评估：价值网络估计无限时域回报
        leaf_value ← V_ψ(node.state)    // ← 关键：替代零终端成本

        // 反向传播：更新Q值（顺序重要！）
        G ← leaf_value
        for node in reversed(path) do
            G ← node.reward + γ · G      // 首先：累积奖励
            node.N ← node.N + 1
            node.W ← node.W + G
            node.Q ← node.W / node.N     // 然后：更新统计量
        end for
    end for

    return argmax_a root.children[a].N
\end{verbatim}

\textbf{PUCT选择公式}： \[
\text{PUCT}(s,a) = Q(s,a) + c \cdot P(a|s) \cdot \frac{\sqrt{N(s)}}{1 + N(s,a)}
\]

其中 \(Q(s,a)\) 是平均回溯价值，\(P(a|s)\)
是来自策略网络的先验（若不可用则为均匀），\(c\) 是探索常数。

\hypertarget{ux4ef7ux503cux7f51ux7edcux4f5cux4e3aadpux7ec8ux7aefux6210ux672c}{%
\paragraph{4.3.1
价值网络作为ADP终端成本}\label{ux4ef7ux503cux7f51ux7edcux4f5cux4e3aadpux7ec8ux7aefux6210ux672c}}

连接SD-MCTS与经典OR的关键洞见是：

\begin{table}[h]
\centering
\begin{tabular}{p{3cm}p{4cm}p{3cm}}
\toprule
\textbf{规划方法} & \textbf{叶节点评估} & \textbf{有效视界} \\
\midrule
MPC（零终端） & \(V_{\text{leaf}} = 0\) & 恰好 \(H\) 步 \\
SD-MCTS & \(V_{\text{leaf}} = V_\psi(s)\) & \(H \times \text{Extension}\) \\
\bottomrule
\end{tabular}
\end{table}

在ADP术语中，\(V_\psi\) 作为\textbf{学习终端成本函数}（Bertsekas
2019），将无限尾部``压缩''为单个标量。这允许短视界搜索考虑长期后果------正是延迟效应系统所需的能力。

\begin{quote}
\textbf{备注（策略网络可选性）}：与AlphaZero中策略网络对于在大动作空间中引导搜索至关重要不同，SD-MCTS在动作空间适中时（如
\(\lvert A\rvert \leq 10\)）可以使用均匀先验
\(P(a\lvert s) = 1/\lvert A\rvert\)
运行。价值网络是关键组件；策略网络主要加速收敛。
\end{quote}

\begin{center}\rule{0.5\linewidth}{0.4pt}\end{center}

\hypertarget{ux8fdeux7eedux72b6ux6001ux7a7aux95f4ux5904ux7406}{%
\subsubsection{4.4
连续状态空间处理}\label{ux8fdeux7eedux72b6ux6001ux7a7aux95f4ux5904ux7406}}

\emph{{[}回应审稿人3关切：``树搜索如何处理连续状态？''{]}}

MCTS传统上为离散状态空间设计。我们如下适配到连续状态：

\textbf{节点定义}：每个节点存储\emph{精确}的连续状态向量
\(s \in \mathbb{R}^n\)。如果两个节点来自根节点的不同动作序列，则它们不同。

\textbf{无状态聚合}：我们不离散化或哈希连续状态。树结构由动作序列决定，而非状态相似性。

\textbf{树复用策略}： - \emph{回合内}：执行动作 \(a^*\) 后，我们复用以
\(\text{root.children}[a^*]\) 为根的子树，避免冗余计算。 -
\emph{跨回合}：由于随机转移和训练期间模型参数变化，树不复用。

\textbf{可扩展性考虑}： - 树深度受搜索视界 \(H\) 限制（通常5-15）。 -
每层树宽度等于 \(|A|\)（对于动作空间适中的OR问题可管理）。 -
每次搜索创建的总节点数：\(O(N_{\text{sim}})\)，与状态维度无关。 -
总模拟操作数：\(O(N_{\text{sim}} \cdot H)\)，每次模拟遍历最多 \(H\)
层（完整复杂度分析见§4.5）。

\begin{quote}
\textbf{何时需要渐进加宽}：对于非常大或连续动作空间的问题（如
\(\lvert A\rvert > 100\)），需要渐进加宽（Coulom
2007）或动作采样。我们当前的测试平台有
\(\lvert A\rvert \leq 5\)，完全扩展是可行的。
\end{quote}

\begin{center}\rule{0.5\linewidth}{0.4pt}\end{center}

\hypertarget{ux8ba1ux7b97ux590dux6742ux5ea6}{%
\subsubsection{4.5 计算复杂度}\label{ux8ba1ux7b97ux590dux6742ux5ea6}}

\hypertarget{ux5728ux7ebfux89c4ux5212ux6210ux672cux6bcfux51b3ux7b56ux6b65}{%
\paragraph{在线规划成本（每决策步）}\label{ux5728ux7ebfux89c4ux5212ux6210ux672cux6bcfux51b3ux7b56ux6b65}}

\[
\text{Cost}_{\text{online}} = O\left(N_{\text{sim}} \times H \times \left(C_{\text{model}} + C_{\text{value}}\right) + C_{\text{tree}}\right)
\]

其中： - \(N_{\text{sim}}\)：MCTS模拟次数（通常20-100） -
\(H\)：搜索深度（通常5-15） -
\(C_{\text{model}}\)：世界模型前向传播（结构化SD：\(O(n)\)，\(n\)
为状态变量数） -
\(C_{\text{value}}\)：价值网络前向传播（MLP：\(O(d_{\text{hidden}}^2)\)）
- \(C_{\text{tree}}\)：树操作（可忽略）

\hypertarget{ux4e0empcux968fux673aux5c04ux51fbux7684ux6bd4ux8f83}{%
\paragraph{与MPC随机射击的比较}\label{ux4e0empcux968fux673aux5c04ux51fbux7684ux6bd4ux8f83}}

\begin{table}[h]
\centering
\begin{tabular}{p{4cm}p{6cm}p{4cm}}
\toprule
\textbf{方法} & \textbf{轨迹采样} & \textbf{随 \(H\) 增长} \\
\midrule
MPC（随机射击） & \(K\) 样本 × \(\lvert A\rvert^H\) 组合 & \textbf{指数级} \(O(K \cdot \lvert A\rvert^H)\) \\
SD-MCTS & \(N_{\text{sim}}\) 模拟 × \(H\) 步 & \textbf{线性} \(O(N_{\text{sim}} \cdot H)\) \\
\bottomrule
\end{tabular}
\end{table}

对于 \(\lvert A\rvert=5\)，\(H=10\)：MPC需要评估 \(\sim 10^7\)
动作序列；SD-MCTS使用 \(\sim 10^3\) 节点评估。

\hypertarget{ux5b9eux8bc1ux8ba1ux65f6ux8be6ux89c17}{%
\paragraph{实证计时（详见§7）}\label{ux5b9eux8bc1ux8ba1ux65f6ux8be6ux89c17}}

\begin{table}[h]
\centering
\begin{tabular}{p{4cm}p{4cm}p{4cm}}
\toprule
\textbf{方法} & \textbf{毫秒/步} & \textbf{相对值} \\
\midrule
SD-MCTS (\(H=1\)) & 57 & 1.0× \\
Oracle MPC (\(H=10\)) & 1812 & 31.8× \\
\bottomrule
\end{tabular}
\end{table}

SD-MCTS以Oracle MPC \textbf{3.2\%} 的计算成本实现相当性能。

\hypertarget{ux5b9eux73b0ux6ce8ux610fux4e8bux9879}{%
\subsubsection{4.6
实现注意事项}\label{ux5b9eux73b0ux6ce8ux610fux4e8bux9879}}

我们记录了几个对可复现性至关重要的实现细节：

\textbf{防止双重计数}：在MCTS反向传播中，叶价值
\(V_\psi(s_{\text{leaf}})\) 应\emph{不}包含叶节点的即时奖励。正确更新是
\(G \leftarrow r + \gamma \cdot V_\psi(s')\)，而非
\(G \leftarrow V_\psi(s')\)（后者会双重计数第一步奖励）。

\textbf{大动作空间的渐进拓宽}：标准MCTS在节点``完全扩展''后才进入selection阶段。对于大动作空间，有限模拟次数永远无法完全扩展，导致每个动作仅被访问一次，Q值无法有效累积。解决方案是引入渐进拓宽：当已扩展子节点数
\(|C| < \sqrt{N}\)（\(N\)为父节点访问次数）时继续扩展新动作，否则在已扩展子节点中用PUCT选择。

\textbf{价值网络的回报归一化}：原始回报范围差异巨大会导致价值网络学习失败，表现为预测与真实回报呈负相关。引入z-score归一化可显著改善学习效果，使相关系数转为强正相关。

\textbf{离散阶跃的梯度传播}：当系统参数涉及离散阶跃变化（如需求在某时刻突变）时，标准硬阈值函数在阈值处梯度为零，阻碍参数学习。我们采用Straight-through
estimator：前向传播使用硬阈值以匹配环境行为，反向传播使用sigmoid软门控以提供梯度信号。

\textbf{不对称损失权重}：在安全关键参数（如效率退化率）学习中，乐观偏差（低估风险）比保守偏差更危险。我们对预测误差施加不对称权重：低估时赋予更高权重，引导模型学习保守估计。

\textbf{参数初始化}：Softplus变换参数需要仔细初始化。对于目标初始值
\(\theta_0\)，我们将原始参数初始化为
\(\phi_0 = \log(\exp(\theta_0) - 1)\)。

\textbf{目标网络稳定性}：价值网络使用软更新的目标网络
\(\bar{\psi} \leftarrow \tau \psi + (1-\tau) \bar{\psi}\)。这防止了当世界模型参数也在变化时TD学习中的振荡。

\begin{center}\rule{0.5\linewidth}{0.4pt}\end{center}

\hypertarget{ux7406ux8bbaux5206ux6790}{%
\subsection{5. 理论分析}\label{ux7406ux8bbaux5206ux6790}}

本节为SD-MCTS提供理论基础。我们呈现三个结果：(i)结构化参数估计的样本复杂度，(ii)通过学习终端值的有效视界扩展，(iii)不完美但深度规划优于完美但浅层规划的条件。

\begin{quote}
\textbf{关于严谨性的说明}：我们明确陈述假设并承认理想化条件（如i.i.d.采样）与算法实践的偏离。这些结果旨在作为阐明机制的\emph{解释性界限}，而非紧致性能保证。
\end{quote}

\hypertarget{ux7ed3ux6784ux5316ux53c2ux6570ux4f30ux8ba1ux7684ux6837ux672cux590dux6742ux5ea6}{%
\subsubsection{5.1
结构化参数估计的样本复杂度}\label{ux7ed3ux6784ux5316ux53c2ux6570ux4f30ux8ba1ux7684ux6837ux672cux590dux6742ux5ea6}}

结构化世界模型的关键优势是降低样本复杂度：学习 \(d\)
个标量参数所需样本远少于学习完整转移表。

\textbf{引理1}（i.i.d.采样下的样本复杂度）

\emph{假设}： 1. 真实动态：\(s' = f(s,a;\theta^*) + \xi\)，其中
\(\theta^* \in \mathbb{R}^d\) 2.
线性参数化：\(f(s,a;\theta) = \Phi(s,a)\theta\)，已知特征
\(\Phi\)，\(\|\Phi(s,a)\|_2 \leq B\) 3.
亚高斯噪声：\(\|\xi\|_{\psi_2} \leq \sigma\) 4.
持续激励：\(\lambda_{\min}(\mathbb{E}[\Phi^\top \Phi]) \geq \lambda_0 > 0\)
5. \textbf{i.i.d.采样}：\(\{(s_i, a_i, \xi_i)\}_{i=1}^N\) 独立

\emph{结论}：以至少 \(1-\delta\) 的概率，最小二乘估计 \(\hat{\theta}_N\)
满足：
\[\|\hat{\theta}_N - \theta^*\|_2 \leq O\left(\sqrt{\frac{d \cdot \log(1/\delta)}{N \cdot \lambda_0}}\right)\]

因此，达到 \(\epsilon\)-精确估计需要 \(N = O(d/\epsilon^2)\) 样本。

\emph{证明梗概}：标准最小二乘分析。Gram矩阵因持续激励而集中于其期望；噪声项由亚高斯尾不等式界定。详见附录B.1。

\begin{quote}
\textbf{备注（与实践的差距）}：假设5（i.i.d.采样）在RL中不满足，因为数据是在策略上收集的并呈现时序相关性。将此结果扩展到轨迹数据需要鞅集中或混合时间论证（见如Simchowitz
et
al.~2018）。我们将引理1作为\emph{解释性}结果展示，说明\emph{为何}结构化模型具有样本效率，而非作为我们算法的形式保证。
\end{quote}

\textbf{推论1}（与表格方法的比较）

表格方法估计 \(|\mathcal{S}||\mathcal{A}|\) 个独立转移概率，每个需要
\(O(1/\epsilon^2)\) 样本，总计
\(O(|\mathcal{S}||\mathcal{A}|/\epsilon^2)\)。当
\(d \ll |\mathcal{S}||\mathcal{A}|\)（结构化假设）时，样本效率增益为
\(|\mathcal{S}||\mathcal{A}|/d\)------可能达到数量级差异。

\begin{center}\rule{0.5\linewidth}{0.4pt}\end{center}

\hypertarget{ux901aux8fc7ux5b66ux4e60ux7ec8ux7aefux503cux7684ux6709ux6548ux89c6ux754cux6269ux5c55}{%
\subsubsection{5.2
通过学习终端值的有效视界扩展}\label{ux901aux8fc7ux5b66ux4e60ux7ec8ux7aefux503cux7684ux6709ux6548ux89c6ux754cux6269ux5c55}}

SD-MCTS的核心机制是使用学习价值网络将有效规划视界扩展到显式搜索深度之外。

\textbf{定理1}（有效视界扩展界）

\emph{假设}： 1.
价值网络误差：\(\|V_\psi - V^*\|_\infty \leq \epsilon_V\) 2.
奖励Lipschitz：\(|R(s,a) - R(s',a)| \leq L_r \|s - s'\|\) 3.
价值Lipschitz：\(|V^*(s) - V^*(s')| \leq L_V \|s - s'\|\) 4.
动态Lipschitz：\(\|f(s,a;\theta) - f(s',a;\theta)\| \leq L_f \|s - s'\|\)
5.
模型误差：\(\|f(s,a;\hat{\theta}) - f(s,a;\theta^*)\| \leq \epsilon_{\text{model}}\)
6. 搜索深度：\(H\)

\emph{结论}：带学习终端值的MCTS策略价值估计误差满足：
\[|V^*(s) - \hat{V}_{\text{MCTS}}(s)| \leq \underbrace{\gamma^H \epsilon_V}_{\text{折扣价值误差}} + \underbrace{(L_r + \gamma^H L_V) \cdot \epsilon_{\text{model}} \cdot \Gamma_H}_{\text{累积模型误差}}\]

其中 \(\Gamma_H = (L_f^H - 1)/(L_f - 1)\)（\(L_f > 1\) 时），或
\(\Gamma_H = H\)（\(L_f \leq 1\) 时）。

\emph{证明}：见附录B.2。关键步骤：(i)将 \(V^*\) 分解为 \(H\)
步奖励加终端值；(ii)通过Lipschitz连续性界定奖励预测误差；(iii)将终端值误差界定为
\(\epsilon_V\) 加状态预测误差传播。

\textbf{推论2}（与MPC的比较）

零终端成本的MPC等价于设置
\(\epsilon_V = \|V^*\|_\infty \leq R_{\max}/(1-\gamma)\)。因此： -
MPC截断误差：\(\gamma^H R_{\max}/(1-\gamma)\) - MCTS+\(V_\psi\)
截断误差：\(\gamma^H \epsilon_V\)

当 \(\epsilon_V \ll R_{\max}/(1-\gamma)\) 时，MCTS+\(V_\psi\)
实现显著更低的截断偏差。

\textbf{实证验证}（见§7.2）：\(H=1\)
的SD-MCTS在延迟崩溃任务上达到100\%成功率，而 \(V=0\) 的MPC需要
\(H \geq 8\)。这对应于\textbf{至少7步的有效视界扩展}。

\begin{quote}
\textbf{备注（有限模拟预算）}：定理1假设MCTS在搜索树内找到最优动作。在有限模拟
\(N_{\text{sim}}\)
下，存在来自不完整树探索的额外误差。根据UCT收敛结果（Kocsis and
Szepesvári 2006），此误差以 \(O(1/\sqrt{N_{\text{sim}}})\)
减少。我们的实验（§7.2，表4）表明SD-MCTS对模拟预算鲁棒：在所有测试的
\(N_{\text{sim}}\)
值下性能保持100\%，而无价值网络的MCTS在低预算下显著退化。
\end{quote}

\begin{center}\rule{0.5\linewidth}{0.4pt}\end{center}

\hypertarget{ux89c6ux754c-ux7cbeux5ea6ux6743ux8861}{%
\subsubsection{5.3
视界-精度权衡}\label{ux89c6ux754c-ux7cbeux5ea6ux6743ux8861}}

我们的核心理论洞见是，在长时间常数系统中，规划深度比参数精度更有价值。

\textbf{定理2}（视界-精度权衡）

\emph{设定}：比较两个规划器： 1. \textbf{Oracle MPC}：完美参数
\(\theta^*\)，视界 \(H\)，零终端成本 2. \textbf{SD-MCTS}：学习参数
\(\hat{\theta}\)（\(\|\hat{\theta} - \theta^*\| \leq \delta_\theta\)），视界
\(H\)，学习 \(V_\psi\)（\(\|V_\psi - V^*\|_\infty \leq \epsilon_V\)）

\emph{结论}：当以下条件满足时，SD-MCTS比Oracle
MPC实现更低的策略价值误差：
\[\underbrace{\frac{\gamma^H R_{\max}}{1-\gamma}}_{\text{MPC截断偏差}} > \underbrace{L_\theta \delta_\theta + \gamma^H \epsilon_V}_{\text{SD-MCTS总误差}}\]

其中 \(L_\theta\) 是 \(V^*\) 关于模型参数的Lipschitz常数。

\emph{解释}：左边是MPC的截断偏差（零终端成本下不可避免）。右边是SD-MCTS来自参数不完美的误差（\(L_\theta \delta_\theta\)）加折扣价值网络误差（\(\gamma^H \epsilon_V\)）。当系统时间常数
\(\tau\) 超过MPC视界 \(H\)
时，截断偏差占主导，SD-MCTS尽管参数不完美仍获胜。

\textbf{数值示例}（\(\gamma = 0.99\)，\(H = 15\)，\(R_{\max} = 10\)）：
-
MPC截断偏差：\(\gamma^{15} \cdot R_{\max}/(1-\gamma) \approx 0.86 \times 1000 = 860\)
- SD-MCTS误差（假设
\(\delta_\theta = 0.25\)，\(L_\theta = 10\)，\(\epsilon_V = 5\)）：\(10 \times 0.25 + 0.86 \times 5 = 6.8\)
- \textbf{比率}：\(860 / 6.8 \approx 126\times\)

截断偏差超过参数诱导误差两个数量级，解释了为何``不完美但远见''胜过``完美但短视''。

\hypertarget{ux76f8ux56fesd-mctsux4f55ux65f6ux5360ux4f18}{%
\paragraph{5.3.1
相图：SD-MCTS何时占优？}\label{ux76f8ux56fesd-mctsux4f55ux65f6ux5360ux4f18}}

我们可以在（视界比，误差比）空间中将权衡条件可视化为\textbf{相图}：

\textbf{轴}： - 横轴：\(H / \tau\)（规划视界 / 系统时间常数） -
纵轴：\(\delta_\theta / \epsilon_V\)（参数误差 / 价值误差）

\textbf{区域}： - \textbf{SD-MCTS占优}：低
\(H/\tau\)（系统相对规划视界有长延迟） - \textbf{Oracle MPC占优}：高
\(H/\tau\)（规划视界覆盖相关动态）

\begin{quote}
【图占位】相图显示SD-MCTS优势区域。§7实验的经验边界叠加在理论预测上。
\end{quote}

\textbf{实践启示}：对于 \(\tau > H\)
的系统（常见于疲劳、磨损、季节性需求），投资于价值网络训练比投资于参数校准产生更高回报。

\hypertarget{ux5de5ux4f5cux793aux4f8bux5ef6ux8fdfux5d29ux6e83ux7684ux4e34ux754cux89c6ux754c}{%
\paragraph{5.3.2
工作示例：延迟崩溃的临界视界}\label{ux5de5ux4f5cux793aux4f8bux5ef6ux8fdfux5d29ux6e83ux7684ux4e34ux754cux89c6ux754c}}

我们为场景2（§6.2.2）解析推导临界视界
\(H_{\text{crit}}\)。该场景描述一个单工人系统：工人累积疲劳会导致不可逆崩溃，同时未来存在一个高价值需求涌动事件，最优策略需要在短期收益与长期风险之间权衡。

\textbf{设定}：疲劳动态 \(f_{t+1} = f_t + 1.1 a_t - 0.1\)，崩溃阈值
\(F = 10\)，涌动在 \(t = 12\)。

\textbf{贪婪轨迹}：\(a_t = 1\) 时，疲劳演化为
\(f_t = t\)（线性增长）。崩溃发生在 \(t^* = 10\)。

\textbf{MPC可见性条件}：对于视界 \(H\)
的MPC能``看到''崩溃并修改行为，我们需要：
\[H \geq t^* - t_{\text{decision}} = 10 - 3 = 7\] 其中
\(t_{\text{decision}} \approx 3\) 是强度降低可以防止崩溃的最后决策点。

\textbf{纳入折扣}：\(\gamma = 0.99\) 时，崩溃惩罚 \(P\) 在视界 \(H\)
处的折扣可见性是 \(\gamma^H P\)。对于MPC优先避免崩溃而非即时奖励
\(r\)，我们需要：
\[\gamma^H P > r \cdot H \implies H > \frac{\log(r \cdot H / P)}{\log(1/\gamma)}\]

数值求解 \(r = 1\)，\(P = 100\)：\(H_{\text{crit}} \approx 6.36\)。

\textbf{实证验证}：表3显示MPC在 \(H \leq 6\) 时达到0\%成功率，\(H = 8\)
时90\%，与 \(H_{\text{crit}} \approx 6-7\) 一致。SD-MCTS在 \(H = 1\)
达到100\%成功率，因为 \(V_\psi\) 隐式编码了崩溃惩罚。

\begin{center}\rule{0.5\linewidth}{0.4pt}\end{center}

\hypertarget{ux7ed3ux6784ux5148ux9a8cux5f15ux5bfcux4e0bux7684ux957fux671fux89c4ux5212ux8befux5deeux53d7ux63a7ux6027}{%
\subsubsection{5.4
结构先验引导下的长期规划误差受控性}\label{ux7ed3ux6784ux5148ux9a8cux5f15ux5bfcux4e0bux7684ux957fux671fux89c4ux5212ux8befux5deeux53d7ux63a7ux6027}}

在运筹学应用中，模型长期预测误差随规划周期的无界累积是极度危险的。SD-MCTS体系通过引入存量-流量领域先验（硬约束）不仅仅是为了降低样本复杂度，更从根本上改变了多步误差的演变行为。

\textbf{定理3}（结构先验下的多步误差上限压缩）

\emph{设定}：考虑具有物理约束限度（如非负库存界 \(s \geq 0\)
和系统瓶颈上线 \(s \leq C\)）的动态系统。若SD-MCTS中的结构化模型转移机制
\(f(s,a;\hat{\theta})\) 注入了状态截流算子 \(\text{clamp}(\cdot, 0, C)\)
作为状态守恒强制约束。

\emph{结论}：多步自回归推演产生的累积状态预测误差将落入不变包络集中，并被压缩至常量界限，不再随视界
\(H\) 爆发，即：
\[ \lim_{H \to \infty} \| \hat{s}_H - s^*_H \| \leq \mathcal{O}(1) \]

\emph{管理学含义}：规划视界的深层延展不会因复用黑盒模型引发长程偏差失控爆炸。这意味着含有守恒结构先验的该体系天然具备``边界内部自我修正''抗风险能力，其保障了深层延迟效应捕捉的可靠基石。这在物理直觉上也深刻契合了存量系统累计求和自带的低通滤波特性——即\textbf{积分平滑能天然抑制高频随机噪声}。

\begin{center}\rule{0.5\linewidth}{0.4pt}\end{center}

\hypertarget{ux7406ux8bbaux8d21ux732eux603bux7ed3}{%
\subsubsection{5.5
理论贡献总结}\label{ux7406ux8bbaux8d21ux732eux603bux7ed3}}

\begin{table}[h]
\centering
\begin{tabular}{p{2.5cm}p{2.5cm}p{9cm}}
\toprule
\textbf{结果} & \textbf{类型} & \textbf{关键洞见} \\
\midrule
命题1 & 条件 & 黑盒预测模型的多步Q估计存在不可消减的结构性指数放大偏差 \\
引理1 & 解释性 & 结构化模型将OR推演的样本辨识复杂度降低 \(\lvert\mathcal{S}\rvert\lvert\mathcal{A}\rvert/d\) 倍 \\
定理1 & 界限 & 学习终端值将全局视界的截断折损从 \(O(R_{\max}/(1-\gamma))\) 降至逼近 \(O(\epsilon_V)\) \\
定理2 & 条件 & 当延迟显著时 \(\gamma^H R_{\max}/(1-\gamma) > L_\theta \delta_\theta + \gamma^H \epsilon_V\)，容忍偏差深搜比完美近搜更奏效 \\
定理3 & 界限 & 实体守恒先验将多步展开截流进入不变可行包络面，长期规划误差由渐进失控钳制回 \(\mathcal{O}(1)\) 量级 \\
\bottomrule
\end{tabular}
\end{table}

\begin{quote}
\textbf{警告}：这些结果依赖于理想化假设（Lipschitz连续性、有界误差、引理1的i.i.d.采样）。它们旨在解释\emph{为何}
SD-MCTS有效，而非提供紧致保证。§7的实证验证确认机制按预测运作。
\end{quote}

\begin{center}\rule{0.5\linewidth}{0.4pt}\end{center}

\hypertarget{ux8ba1ux7b97ux7814ux7a76}{%
\subsection{6. 计算研究}\label{ux8ba1ux7b97ux7814ux7a76}}

本节描述我们的实验设计、测试平台和基准方法。我们遵循IJOC计算实验指南，强调可复现性、公平比较和统计严谨性。

\hypertarget{ux5b9eux9a8cux8bbeux8ba1ux539fux5219}{%
\subsubsection{6.1
实验设计原则}\label{ux5b9eux9a8cux8bbeux8ba1ux539fux5219}}

\hypertarget{ux516cux5e73ux6027ux51c6ux5219}{%
\paragraph{6.1.1 公平性准则}\label{ux516cux5e73ux6027ux51c6ux5219}}

我们对所有比较执行两条公平性准则：

\textbf{样本预算对齐}：所有基于学习的方法在训练期间接收相同数量的环境交互。
- 训练预算：30,000环境步 - 约100回合，每回合300步

\textbf{计算预算对齐}：我们报告在线决策时间（毫秒/步）以实现性能/计算比较。
- 所有方法在相同硬件上运行（见§8） - 我们变化规划参数（视界
\(H\)，模拟次数 \(N_{\text{sim}}\)）以追踪Pareto前沿

\hypertarget{ux8bc4ux4f30ux6307ux6807}{%
\paragraph{6.1.2 评估指标}\label{ux8bc4ux4f30ux6307ux6807}}

\begin{table}[h]
\centering
\begin{tabular}{p{5cm}p{6cm}p{4cm}}
\toprule
\textbf{指标} & \textbf{定义} & \textbf{目的} \\
\midrule
平均回报 & \(\frac{1}{N}\sum_{i=1}^N \sum_{t} \gamma^t r_t^{(i)}\) & 主要性能度量 \\
成功率 & 避免失败的回合比例 & 对安全敏感任务关键 \\
崩溃/失败率 & 以系统失败结束的回合比例 & 捕获灾难性结果 \\
涌动捕获率 & 成功捕获高价值事件的比例 & 衡量长视界规划 \\
约束违反率 & 违反物理约束的状态比例 & 测试模型有效性 \\
在线时间 & 每决策的挂钟时间 & 计算效率 \\
\bottomrule
\end{tabular}
\end{table}

\hypertarget{ux7edfux8ba1ux62a5ux544a}{%
\paragraph{6.1.3 统计报告}\label{ux7edfux8ba1ux62a5ux544a}}

\begin{itemize}
\tightlist
\item
  \textbf{随机种子}：根据各场景的统计特性需要，使用10至50个随机种子
\item
  \textbf{不确定性度量}：表格中``±''后数值为标准差(SD)；统计检验时报告95\%
  CI
\item
  \textbf{显著性检验}：配对t检验，多重比较使用Bonferroni校正
\item
  \textbf{效应量}：Cohen's \(d\) 用于实际显著性
\end{itemize}

\begin{center}\rule{0.5\linewidth}{0.4pt}\end{center}

\hypertarget{ux6d4bux8bd5ux5e73ux53f0}{%
\subsubsection{6.2 测试平台}\label{ux6d4bux8bd5ux5e73ux53f0}}

我们在三个测试平台上评估，设计用于隔离关键机制：(i)渐进效率退化，(ii)带未来高价值事件的延迟硬崩溃，以及(iii)多级供应链中的牛鞭效应。

\hypertarget{ux573aux666f1ux5e26ux6548ux7387ux9000ux5316ux7684ux4efbux52a1ux5206ux914d}{%
\paragraph{6.2.1
场景1：带效率退化的任务分配}\label{ux573aux666f1ux5e26ux6548ux7387ux9000ux5316ux7684ux4efbux52a1ux5206ux914d}}

\textbf{环境}：\(N=2\) 工人处理随机到达的任务分配系统。

\textbf{动态}： - 到达：\(A_t \sim \text{Poisson}(\lambda \Delta t)\) -
服务率：\(s_{i,t} = \eta(f_t) \cdot \mu_i\)，效率随疲劳退化 -
效率函数：\(\eta(f) = 1/(1 + \alpha \cdot \max(0, f-\tau)^2)\)

\textbf{参数}：

\begin{table}[h]
\centering
\begin{tabular}{p{3cm}p{3cm}p{5cm}p{3cm}}
\toprule
\textbf{参数} & \textbf{真值} & \textbf{初始信念} & \textbf{误差} \\
\midrule
\(\alpha\)（退化率） & 0.1 & 0.001 & 100× \\
\(\tau\)（疲劳阈值） & 5.0 & 12.0 & 2.4× \\
\bottomrule
\end{tabular}
\end{table}

\textbf{奖励}：负队列长度和积压（最小化等待）

\textbf{挑战}：疲劳累积的延迟效应在短规划视界内不可见。\(H < 10\)
的MPC系统性地过度使用工人，导致效率崩溃。

\textbf{OR相关性}：此场景代表劳动力调度、机器利用和任何具有性能退化动态的系统。

\hypertarget{ux573aux666f2ux5e26ux9700ux6c42ux6d8cux52a8ux7684ux5ef6ux8fdfux5d29ux6e83}{%
\paragraph{6.2.2
场景2：带需求涌动的延迟崩溃}\label{ux573aux666f2ux5e26ux9700ux6c42ux6d8cux52a8ux7684ux5ef6ux8fdfux5d29ux6e83}}

\textbf{环境}：具有硬崩溃阈值和已知未来高价值事件的单工人系统。

\textbf{动态}： \[f_{t+1} = f_t + (\kappa + \beta) \cdot a_t - \beta\] -
\(f_t\)：累积疲劳 -
\(a_t \in \{0, 0.25, 0.5, 0.75, 1.0\}\)：工作强度（离散） -
\(\kappa = 1.0\)：疲劳累积率 - \(\beta = 0.1\)：恢复率

\textbf{崩溃机制}：若
\(f_t \geq F\)，系统进入吸收失败状态（零未来奖励）。

\textbf{需求涌动}：在 \(t = T_s = 12\) 时，若
\(f_t \leq F_{\text{req}} = 5\)，智能体获得奖金 \(B = 15\)。

\textbf{参数}：

\begin{table}[h]
\centering
\begin{tabular}{p{3.5cm}p{2.5cm}p{2.5cm}p{4cm}}
\toprule
\textbf{参数} & \textbf{真值} & \textbf{MPC信念} & \textbf{误差} \\
\midrule
\(F\)（崩溃阈值） & 10 & 50 & 5×高估 \\
\bottomrule
\end{tabular}
\end{table}

\textbf{挑战}：贪婪策略（\(a_t = 1\)）最大化短期奖励但在
\(t \approx 10\)
时导致崩溃，错过涌动奖金。最优策略需要克制：\(a^* \approx 0.47\)。

\textbf{OR相关性}：代表维护调度、安全关键运营和任何具有灾难性失败模式和时间敏感机会的系统。

\hypertarget{ux573aux666f3ux5564ux9152ux6e38ux620fux4f9bux5e94ux94feux725bux97adux6548ux5e94}{%
\paragraph{6.2.3
场景3：啤酒游戏（供应链牛鞭效应）}\label{ux573aux666f3ux5564ux9152ux6e38ux620fux4f9bux5e94ux94feux725bux97adux6548ux5e94}}

\textbf{环境}：经典MIT啤酒游戏（Sterman
1989），一个多级供应链模拟，展示牛鞭效应------需求变异性随订单向上游传播而放大。

\textbf{结构}：四级供应链：零售商 → 批发商 → 分销商 → 工厂

\begin{verbatim}
客户 → [零售商] → [批发商] → [分销商] → [工厂]
          ↓           ↓          ↓          ↓
        库存        库存        库存        库存
          ↑           ↑          ↑          ↑
        订单        订单        订单        生产
\end{verbatim}

\textbf{存量-流量动态}（每级 \(i\)）：
\[I_{i,t+1} = I_{i,t} + \text{到货}_{i,t} - \text{发货}_{i,t}\]
\[B_{i,t+1} = B_{i,t} + \text{订单}_{i,t} - \text{发货}_{i,t}\]
\[\text{在途}_{i,t+1} = \text{在途}_{i,t} + \text{订单}_{i,t} - \text{到货}_{i,t}\]

其中： - \(I_{i,t}\)：级 \(i\) 的现有库存 -
\(B_{i,t}\)：积压（未满足订单） -
\(\text{在途}_{i,t}\)：已下但未收到的订单 -
信息延迟：\(\tau_{\text{info}} = 2\) 期 -
运输延迟：\(\tau_{\text{ship}} = 2\) 期

\textbf{动作空间}：每级的订单数量
\(o_{i,t} \in \{0, 1, 2, \ldots, O_{\max}\}\)（离散）

\textbf{奖励}：所有级别的负总成本：
\[R_t = -\sum_{i=1}^{4} \left(h \cdot I_{i,t}^+ + p \cdot B_{i,t}\right)\]
其中
\(h = 0.5\)（持有成本），\(p = 1.0\)（积压惩罚），\(I^+ = \max(0, I)\)。

\textbf{需求过程}：客户需求遵循阶跃函数（经典啤酒游戏）： -
\(D_t = 4\)，当 \(t < 5\) - \(D_t = 8\)，当 \(t \geq 5\)（永久阶跃增加）

\textbf{可学习参数}： - 提前期 \(\tau_i\)（若视为不确定） - 需求分布参数
- 成本系数（用于敏感性分析）

\textbf{实验设计}：我们设计两个互补实验以验证SD-MCTS在供应链场景下的优势：

\begin{enumerate}
\def\labelenumi{\arabic{enumi}.}
\tightlist
\item
  \textbf{实验3A：视界墙效应（Horizon Wall）}

  \begin{itemize}
  \tightlist
  \item
    目标：验证价值网络 \(V_\psi\) 的有效视界扩展能力
  \item
    设置：标准啤酒游戏（\(T=36\)），30个评估回合
  \item
    评估指标：总成本、牛鞭比率、计算时间
  \end{itemize}
\item
  \textbf{实验3B：供应链韧性（Resilience）}

  \begin{itemize}
  \tightlist
  \item
    目标：测试面对需求冲击的适应能力
  \item
    设置：需求在 \(t=15\)
    从4突增至16（4倍冲击），持续15期后恢复（\(T=60\)）
  \item
    评估指标：峰值积压、恢复时间、累积超额成本
  \end{itemize}
\item
  \textbf{实验3C：稳健性与管理情景压力测试（Causal Intervention Stress
  Tests）}

  \begin{itemize}
  \tightlist
  \item
    目标：测试决策模型在缺乏重训条件下应对极端结构性突变（如工厂意外停工、物流通信断链）的策略韧性。
  \item
    设置：在推理期对预训练环境的基础底层拓扑实施人为不可导干预（强行斩断下达信息流，或物理截断容量）。
  \item
    评估指标：系统服务率跌幅、策略崩溃失控率、超额成本涨幅。
  \end{itemize}
\end{enumerate}

\textbf{基线方法}：

\begin{table}[h]
\centering
\begin{tabular}{p{4cm}p{6.5cm}p{3.5cm}}
\toprule
\textbf{方法} & \textbf{描述} & \textbf{实验} \\
\midrule
\textbf{Base-Stock} & \((s, S)\) 库存策略，基于目标库存水平订购 & 3A, 3B \\
\textbf{Anchoring-Adjustment} & Sterman (1989) 行为模型，模拟人类启发式决策 & 3A \\
\textbf{Constant-4} & 恒定订购策略（每期订购4单位） & 3A \\
\textbf{Reactive} & 反应式策略（订单 = 观测需求） & 3A \\
\textbf{MPC-H\(n\)} & \(n\) 步前瞻MPC（\(n \in \{2, 4, 6, 8, 10, 12, 16, 20\}\)） & 3A, 3B \\
\textbf{SD-MCTS} & 我们的方法：SD世界模型 + \(V_\psi\) & 3A, 3B \\
\bottomrule
\end{tabular}
\end{table}

\textbf{OR相关性}：啤酒游戏是供应链管理中系统动力学的典型演示。牛鞭效应（Lee
et al.~1997）和供应链韧性（Christopher \& Peck
2004）是当前OR研究的核心议题。

\hypertarget{ux89c4ux6a21ux6269ux5c55-ux5f85ux8865ux5145}{%
\paragraph{6.2.4 规模扩展
{[}待补充{]}}\label{ux89c4ux6a21ux6269ux5c55-ux5f85ux8865ux5145}}

为回应审稿人对问题规模的关切，我们计划以下扩展：

\begin{table}[h]
\centering
\begin{tabular}{p{3cm}p{6cm}p{5cm}}
\toprule
\textbf{扩展} & \textbf{描述} & \textbf{目的} \\
\midrule
\textbf{N-scaling} & 将工人从 \(N=2\) 增加到 \(N=5, 10, 20\) & 测试MCTS随动作空间的可扩展性 \\
\textbf{参数漂移} & 在测试时将 \(\alpha, \tau\) 偏移±30\% & 测试对分布偏移的鲁棒性 \\
\textbf{随机到达} & 比较平均场与真实泊松采样 & 验证平均场近似 \\
\bottomrule
\end{tabular}
\end{table}

\begin{center}\rule{0.5\linewidth}{0.4pt}\end{center}

\hypertarget{ux6bd4ux8f83ux65b9ux6cd5}{%
\subsubsection{6.3 比较方法}\label{ux6bd4ux8f83ux65b9ux6cd5}}

\hypertarget{ux73b0ux6709ux57faux7ebf}{%
\paragraph{6.3.1 现有基线}\label{ux73b0ux6709ux57faux7ebf}}

\begin{table}[h]
\centering
\begin{tabular}{p{4cm}p{2.5cm}p{2.5cm}p{2.5cm}p{2.5cm}}
\toprule
\textbf{方法} & \textbf{参数} & \textbf{视界} & \textbf{终端值} & \textbf{学习} \\
\midrule
MPC（错误信念） & 初始错误 \(\hat{\theta}\) & \(H=3\) & \(V=0\) & 无 \\
自适应MPC & 在线RLS更新 & \(H=3\)或\(H=15\) & \(V=0\) & RLS \\
Oracle MPC & 真实 \(\theta^*\) & \(H=15\) & \(V=0\) & 无（oracle） \\
SD-MCTS & 学习 \(\hat{\theta}\) & \(H\) + \(V_\psi\) & 学习 & 神经网络 \\
\bottomrule
\end{tabular}
\end{table}

\textbf{聚焦MPC的理由}：我们主要与MPC变体比较因为： 1.
MPC是OR中动态优化的主导范式 2. 比较隔离我们的核心贡献：学习终端值 3.
无模型RL（PPO、SAC）需要 \(\sim 10^6\) 样本，使样本效率比较不公平

\hypertarget{ux6837ux672cux6548ux7387ux57faux7ebfux6bd4ux8f83}{%
\paragraph{6.3.2
样本效率基线比较}\label{ux6837ux672cux6548ux7387ux57faux7ebfux6bd4ux8f83}}

为验证结构化世界模型的样本效率优势，我们将SD-MCTS与以下基线进行对比：

\begin{table}[h]
\centering
\begin{tabular}{p{3cm}p{5cm}p{6cm}}
\toprule
\textbf{方法} & \textbf{描述} & \textbf{世界模型类型} \\
\midrule
\textbf{SD-MCTS} & 本文方法 & 结构化（系统动力学） \\
\textbf{BlackBox-MCTS} & MLP动态 + MCTS & 神经网络（无结构） \\
\textbf{DreamerV3} & 端到端世界模型RL & 神经网络（RSSM） \\
\bottomrule
\end{tabular}
\end{table}

\textbf{实验设置}： - 环境：啤酒游戏供应链（Beer
Game），4层级，每回合36步 - 样本预算：1000步（极低样本场景） -
评估：10个固定种子回合的平均回报 -
动作空间：81个离散动作（3\^{}4，每层级3个订货量选择）

\textbf{运行指令}：

\begin{verbatim}
cd docs/neurips && python -u experiments/run_comparison.py \
    --budgets 1000 --seeds 1 --eval-episodes 10 --mcts-sims 100
\end{verbatim}

\textbf{实验结果}（日志位置：\texttt{docs/neurips/results/comparison/20260203\_095159/}）：

\begin{table}[h]
\centering
\begin{tabular}{p{3.5cm}p{6cm}p{4cm}}
\toprule
\textbf{方法} & \textbf{回报（均值±标准差）} & \textbf{训练时间} \\
\midrule
\textbf{SD-MCTS} & \textbf{-649.6 ± 117.1} & 29.0s \\
BlackBox-MCTS & -788.8 ± 101.8 & 76.0s \\
DreamerV3 & -1697.6 ± 127.9 & 49.0s \\
\bottomrule
\end{tabular}
\end{table}

\textbf{分析}： 1.
\textbf{SD-MCTS显著优于所有基线}：在仅1000步的极低样本预算下，SD-MCTS实现了最佳回报（-649.6），比BlackBox-MCTS提升17\%，比DreamerV3提升62\%。

\begin{enumerate}
\def\labelenumi{\arabic{enumi}.}
\setcounter{enumi}{1}
\item
  \textbf{结构先验的价值}：BlackBox-MCTS使用相同的MCTS规划但无结构化世界模型，性能显著下降，验证了存量-流量结构的归纳偏置贡献。
\item
  \textbf{端到端学习的局限}：DreamerV3作为前沿的神经世界模型方法，在低样本场景下表现最差，原因是其需要从零学习整个动态模型，而1000步数据不足以形成有效的世界表示。
\item
  \textbf{训练效率}：SD-MCTS不仅样本效率最高，训练时间也最短（29s），因为结构化模型的参数空间远小于神经网络。
\end{enumerate}

\begin{quote}
\textbf{关键发现}：在OR典型的数据稀缺场景（\(N \sim 10^3\)）下，领域结构知识可将样本效率提升2.6倍以上。这验证了引理1的理论分析：结构化模型的样本复杂度为
\(O(d/\epsilon^2)\)，远低于黑盒模型的
\(O(|\mathcal{S}||\mathcal{A}|/\epsilon^2)\)。
\end{quote}

\hypertarget{ux4e3aux4ec0ux4e48ux4e0dux6bd4ux8f83ux65e0ux6a21ux578brl}{%
\paragraph{6.3.3
为什么不比较无模型RL？}\label{ux4e3aux4ec0ux4e48ux4e0dux6bd4ux8f83ux65e0ux6a21ux578brl}}

我们不将无模型基线（PPO、SAC、DQN）作为主要比较因为：

\begin{enumerate}
\def\labelenumi{\arabic{enumi}.}
\tightlist
\item
  \textbf{样本效率不匹配}：无模型方法通常需要 \(10^6\)+
  样本；我们的预算是 \(3 \times 10^4\)
\item
  \textbf{假设不同}：无模型从零学习；我们假设结构知识
\item
  \textbf{研究焦点}：我们的贡献是\emph{利用结构}，而非超越无模型RL
\end{enumerate}

样本效率比较（回报 vs.~训练步数）可能包含在补充材料中。

\begin{center}\rule{0.5\linewidth}{0.4pt}\end{center}

\hypertarget{ux5b9eux73b0ux7ec6ux8282}{%
\subsubsection{6.4 实现细节}\label{ux5b9eux73b0ux7ec6ux8282}}

\textbf{世界模型}： - 通过softplus变换初始化参数 -
学习率：\(\eta_\theta = 0.01\)，优化器：Adam

\textbf{价值网络}： - 架构：MLP，2隐层 × 64单元，ReLU -
学习率：\(\eta_\psi = 0.001\) - 软更新目标网络（\(\tau = 0.01\)） -
正则化：L2权重衰减 \(\lambda = 10^{-4}\)，无dropout（小网络） -
早停：若验证TD误差连续10回合上升则终止训练 -
训练/验证分割：回放缓冲区的80\%/20\%，每回合刷新

\textbf{MCTS}： - 模拟次数：\(N_{\text{sim}} = 50\)（默认） -
探索常数：\(c = 1.0\) - 折扣因子：\(\gamma = 0.99\)

\textbf{训练}：100回合 × 300步，回放缓冲区10,000，批大小64。

完整超参数见附录D。

\begin{center}\rule{0.5\linewidth}{0.4pt}\end{center}

\hypertarget{ux7ed3ux679c}{%
\subsection{7. 结果}\label{ux7ed3ux679c}}

我们按以下组织呈现实验结果：(i)主要性能比较，(ii)隔离关键组件的机制研究，(iii)可扩展性和鲁棒性分析，(iv)实践洞见。

\hypertarget{ux4e3bux8981ux6027ux80fdux6bd4ux8f83}{%
\subsubsection{7.1
主要性能比较}\label{ux4e3bux8981ux6027ux80fdux6bd4ux8f83}}

\hypertarget{ux573aux666f1ux5e26ux6548ux7387ux9000ux5316ux7684ux4efbux52a1ux5206ux914d-1}{%
\paragraph{7.1.1
场景1：带效率退化的任务分配}\label{ux573aux666f1ux5e26ux6548ux7387ux9000ux5316ux7684ux4efbux52a1ux5206ux914d-1}}

\textbf{表1：主要结果 --- 场景1（50个种子，30,000训练步）}

\begin{table}[h]
\centering
\begin{tabular}{p{3.5cm}p{3.5cm}p{2.5cm}p{2.5cm}p{2cm}}
\toprule
\textbf{方法} & \textbf{平均回报} & \textbf{vs. 基线} & \textbf{Cohen's \(d\)} & \textbf{毫秒/步} \\
\midrule
MPC（错误信念） & −8469 ± 506 & --- & --- & 12 \\
自适应MPC (H=15) & −8154 ± 476 & +3.7\% & 0.64 & 18 \\
Oracle MPC (H=15) & −8098 ± 471 & +4.4\% & 0.75 & 15 \\
\textbf{SD-MCTS} & \textbf{−7853 ± 385} & \textbf{+7.3\%} & \textbf{2.05} & 57 \\
\bottomrule
\end{tabular}
\end{table}

\textbf{关键发现}：SD-MCTS显著超越Oracle MPC
\textbf{3.0\%}（配对t检验：\(t = 14.52\)，\(p < 0.0001\)，Cohen's
\(d = 2.05\)）。

这是一个惊人的结果：\emph{学习}参数的SD-MCTS击败了\emph{完美}参数的Oracle
MPC。优势源于价值网络能够预测MPC规划视界之外的长期效率退化。

\textbf{参数恢复}：30,000步后，SD-MCTS恢复： - \(\alpha\)：0.001 →
0.075（真值：0.1，恢复75\%） - \(\tau\)：12.0 →
4.72（真值：5.0，恢复94\%）

\begin{quote}
【图占位】图1：带95\% CI的性能比较柱状图 【图占位】图2：学习效率函数
vs.~真实 vs.~初始信念
\end{quote}

\hypertarget{ux573aux666f2ux5e26ux9700ux6c42ux6d8cux52a8ux7684ux5ef6ux8fdfux5d29ux6e83-1}{%
\paragraph{7.1.2
场景2：带需求涌动的延迟崩溃}\label{ux573aux666f2ux5e26ux9700ux6c42ux6d8cux52a8ux7684ux5ef6ux8fdfux5d29ux6e83-1}}

\textbf{表2：主要结果 --- 场景2（10个种子）}

\begin{table}[h]
\centering
\begin{tabular}{p{4cm}p{3.5cm}p{2.5cm}p{3.5cm}}
\toprule
\textbf{方法} & \textbf{平均回报} & \textbf{崩溃率} & \textbf{涌动捕获} \\
\midrule
MPC（错误信念） & 9.0 ± 0.0 & \textbf{100\%} & 0\% \\
\textbf{SD-MCTS} & \textbf{15.0 ± 0.0} & \textbf{0\%} & \textbf{100\%} \\
最优（参考）† & 24.8 & 0\% & 100\% \\
\bottomrule
\end{tabular}
\end{table}

†最优通过完美信息下的动态规划计算。

\textbf{关键发现}：SD-MCTS比MPC实现\textbf{66.7\%的改进}。MPC贪婪最大化强度，在
\(t \approx 10\)
崩溃，错过涌动奖金。SD-MCTS学会克制强度，避免崩溃并捕获所有涌动奖金。

\textbf{涌现策略
---``主动休息''}：SD-MCTS发现了反直觉策略：将工作强度降低到
\(a \approx 0.47\)（vs.~贪婪
\(a = 1.0\)），牺牲短期奖励以确保生存和涌动捕获。

\textbf{机制解释}：安全强度阈值可以解析推导。给定疲劳动态
\(f_{t+1} = f_t + (\kappa + \beta) \cdot a_t - \beta\)，恒定强度 \(a\)
下的稳态疲劳是
\(f_\infty = (\kappa + \beta) a / \beta - 1\)。为在涌动时间 \(T_s = 12\)
满足 \(f_{T_s} \leq F_{\text{req}} = 5\)，最大安全强度是：
\[a_{\text{safe}} = \frac{\beta \cdot (F_{\text{req}} + 1)}{(\kappa + \beta) \cdot T_s} \approx 0.45\]
SD-MCTS收敛到
\(a \approx 0.47\)，在解析界限的5\%以内------展示价值网络隐式捕获了长视界约束。

\begin{quote}
【图占位】图3：疲劳轨迹比较（MPC vs.~SD-MCTS）
【图占位】图4：随时间累积奖励
\end{quote}

\hypertarget{ux573aux666f3ux5564ux9152ux6e38ux620fux4f9bux5e94ux94feux52a8ux6001}{%
\paragraph{7.1.3
场景3：啤酒游戏（供应链动态）}\label{ux573aux666f3ux5564ux9152ux6e38ux620fux4f9bux5e94ux94feux52a8ux6001}}

\textbf{实验3A：视界墙效应（Horizon Wall）}

我们首先验证SD-MCTS的有效视界扩展能力。在标准啤酒游戏设置下（\(T=36\),
\(T_s=5\),
\(D_h=8\)），比较SD-MCTS与不同规划视界的MPC。值得注意的是，SD-MCTS需要从初始信念开始学习供应链动态参数（如提前期、需求模式），而MPC基线则假设已知完美参数。

\textbf{表2c：视界墙实验结果（30回合）}

\begin{table}[h]
\centering
\begin{tabular}{p{2.5cm}p{3cm}p{2.5cm}p{2.5cm}p{2.5cm}p{2cm}}
\toprule
\textbf{方法} & \textbf{总成本} & \textbf{95\% CI} & \textbf{vs SD-MCTS} & \textbf{牛鞭比率(R)} & \textbf{时间(s/回合)} \\
\midrule
\textbf{SD-MCTS} & \textbf{623.6 ± 150.0} & [569.9, 677.3] & --- & 17.6 & --- \\
MPC-H2 & 3039.7 ± 1230.0 & [2599.5, 3479.8] & +387.4\% & 22.0 & 0.62 \\
MPC-H4 & 702.0 ± 164.5 & [643.1, 760.9] & +12.6\% & 20.1 & 1.14 \\
MPC-H6 & 665.8 ± 134.7 & [617.6, 714.0] & +6.8\% & 18.9 & 1.61 \\
\textbf{MPC-H8} & \textbf{580.9 ± 64.3} & [557.9, 603.9] & \textbf{−6.8\%} & 19.2 & 1.99 \\
MPC-H10 & 599.2 ± 48.7 & [581.8, 616.6] & −3.9\% & 19.1 & 2.44 \\
MPC-H12 & 642.9 ± 51.6 & [624.4, 661.3] & +3.1\% & 20.0 & 3.02 \\
MPC-H16 & 726.9 ± 59.4 & [705.7, 748.2] & +16.6\% & 21.0 & 4.01 \\
MPC-H20 & 828.3 ± 76.7 & [800.9, 855.8] & +32.8\% & 21.8 & 4.89 \\
\bottomrule
\end{tabular}
\end{table}

†牛鞭比率(R) = Var(零售商订单)/Var(客户需求)，遵循Lee et al.
(1997)定义。表中±后数值为标准差(SD)。

\textbf{关键发现}：

\begin{enumerate}
\def\labelenumi{\arabic{enumi}.}
\item
  \textbf{视界墙现象}：MPC性能呈现非单调模式------H8达到最优（580.9），之后性能\emph{恶化}。这是经典的``视界墙''效应：过长视界导致优化目标偏离当前决策。
\item
  \textbf{SD-MCTS等效视界}：SD-MCTS（623.6）性能介于MPC-H6（665.8）和MPC-H8（580.9）之间，\textbf{等效于约H7的MPC规划能力}。考虑到SD-MCTS的名义树深度约为10，\(V_\psi\)
  提供了约−3步的``负扩展''------这符合预期，因为值网络需要学习适应供应链的多级延迟结构。
\item
  \textbf{统计显著性}：SD-MCTS与MPC-H8差异不显著（\(t=1.39\),
  \(p=0.17\)），两者性能可比。但SD-MCTS显著优于MPC-H2/H4/H16/H20（\(p<0.05\)）。
\end{enumerate}

\textbf{与启发式基线比较}：

\begin{table}[h]
\centering
\begin{tabular}{p{4cm}p{4cm}p{3.5cm}p{3cm}}
\toprule
\textbf{方法} & \textbf{总成本} & \textbf{vs SD-MCTS} & \textbf{牛鞭比率(R)} \\
\midrule
\textbf{SD-MCTS} & \textbf{623.6 ± 150.0} & --- & 17.6 \\
Reactive & 861.2 ± 263.2 & +38.1\% & 3.3 \\
Constant-4 & 2362.4 ± 239.7 & +278.9\% & 0.0 \\
Base-Stock & 2458.4 ± 829.7 & +294.3\% & 12.6 \\
Anchoring-Adjustment & 5423.4 ± 561.7 & +770.0\% & 15.2 \\
\bottomrule
\end{tabular}
\end{table}

SD-MCTS显著优于所有启发式基线，成本降低38\%至770\%。

\begin{center}\rule{0.5\linewidth}{0.4pt}\end{center}

\textbf{实验3B：供应链韧性（Resilience）}

我们测试SD-MCTS面对需求冲击的适应能力。在 \(t=15\)
时需求从4突增至16（4倍冲击），持续至
\(t=30\)。尽管SD-MCTS在训练阶段仅暴露于正常需求模式，其学习的价值网络仍能有效应对分布外的需求涌动。

\textbf{表2d：供应链韧性实验结果（30回合，10000训练步）}

\begin{table}[h]
\centering
\begin{tabular}{p{3.5cm}p{3.5cm}p{2.5cm}p{3cm}p{2cm}}
\toprule
\textbf{方法} & \textbf{总成本} & \textbf{峰值积压} & \textbf{恢复时间(步)} & \textbf{最大退化} \\
\midrule
\textbf{SD-MCTS} & \textbf{2024.1 ± 418.5} & \textbf{68.5 ± 23.3} & 12.4 ± 7.8 & 2.87 \\
MPC-H4 & 2036.3 ± 540.3 & 70.5 ± 30.5 & 11.3 ± 9.0 & 3.27 \\
MPC-H8 & 2295.9 ± 224.6 & 91.5 ± 14.2 & 11.6 ± 3.3 & 3.67 \\
MPC-H12 & 2465.1 ± 153.3 & 94.3 ± 9.4 & 10.6 ± 4.7 & 3.64 \\
Base-Stock & 10358.0 ± 0.0 & 360.0 ± 0.0 & 30.0 ± 0.0 & 12.30 \\
\bottomrule
\end{tabular}
\end{table}

*注：表中±后数值为标准差(SD)。

\textbf{核心发现：最优规划视界随环境动态变化}

对比实验3A（稳态）和3B（冲击）的MPC结果，我们观察到一个重要现象：

\begin{table}[h]
\centering
\begin{tabular}{p{3.5cm}p{3cm}p{2.5cm}p{2.5cm}p{2.5cm}}
\toprule
\textbf{环境条件} & \textbf{MPC最优视界} & \textbf{最优成本} & \textbf{次优视界} & \textbf{次优成本} \\
\midrule
稳态 (3A) & H8 & 580.9 & H6 & 665.8 \\
需求冲击 (3B) & \textbf{H4} & \textbf{2036.3} & H8 & 2295.9 \\
\bottomrule
\end{tabular}
\end{table}

\textbf{为什么冲击环境下短视界更优？}
这一反直觉现象源于两个相互强化的机制：

\begin{quote}
\textbf{机制解释框：H4 \textgreater{} H8 悖论的双重成因}

\begin{enumerate}
\def\labelenumi{\arabic{enumi}.}
\item
  \textbf{预测锚定偏差（Forecast Anchor
  Bias）}：SD模型的隐式稳态假设在长视界回滚中产生``预测锚点''。当需求从4突增至16时，H8优化的是``需求=4''的预期未来，而非``需求=16''的当前现实。较长视界在更多时间步上累积这种偏差。
\item
  \textbf{Jacobian误差放大}：多步回滚误差按
  \(\Delta s_k \approx (\prod_{j} J_f^{(j)}) \cdot \Delta\theta\)
  增长。当转移Jacobian的谱半径 \(> 1\)
  时，参数误差随视界深度指数级放大，导致价值估计不稳定。
\end{enumerate}

\textbf{短视界的``无知之福''}：H4仅对当前积压做出反应，减少了对错误预测的暴露，反而在冲击条件下表现更优。
\end{quote}

\textbf{SD-MCTS的跨环境鲁棒性}：SD-MCTS无需手动选择规划视界，其学习的价值网络在不同环境条件下均能达到接近最优MPC的性能：
- 在稳态下，SD-MCTS (623.6) 性能接近最优 MPC-H8 (580.9) -
在冲击下，SD-MCTS (2024.1) 性能接近最优 MPC-H4 (2036.3)

这种鲁棒性对实践者具有重要价值：在真实供应链中，需求模式可能在稳态和冲击之间切换，手动调优MPC视界既困难又滞后，而SD-MCTS无需重新调优即可在不同环境下保持竞争力。

\textbf{统计显著性分析}：

\begin{table}[h]
\centering
\begin{tabular}{p{3cm}p{3cm}p{3cm}p{1.5cm}p{2cm}p{2cm}}
\toprule
\textbf{比较对} & \textbf{成本差异} & \textbf{95\% CI} & \textbf{p值} & \textbf{Cohen's d} & \textbf{结论} \\
\midrule
SD-MCTS vs MPC-H4 & −12.3 (−0.6\%) & [−262.1, +237.6] & 0.92 & 0.02 & 无显著差异 \\
SD-MCTS vs MPC-H8 & −271.9 (−11.8\%) & [−433.3, −110.4] & 0.0026** & 0.60 & SD-MCTS显著更优 \\
\bottomrule
\end{tabular}
\end{table}

\textbf{峰值积压分析}：

\begin{table}[h]
\centering
\begin{tabular}{p{4cm}p{3.5cm}p{2cm}p{3.5cm}}
\toprule
\textbf{比较对} & \textbf{差异} & \textbf{p值} & \textbf{Cohen's d} \\
\midrule
SD-MCTS vs MPC-H4 & −2.0 (−2.8\%) & 0.78 & 0.07 ns \\
SD-MCTS vs MPC-H8 & −22.9 (−25.1\%) & 2.9e-5*** & \textbf{0.90 大} \\
\bottomrule
\end{tabular}
\end{table}

\textbf{尾部风险分析（90分位）}：

\begin{table}[h]
\centering
\begin{tabular}{p{3.5cm}p{3.5cm}p{3.5cm}p{3.5cm}}
\toprule
\textbf{方法} & \textbf{成本 P90} & \textbf{峰值积压 P90} & \textbf{恢复时间 P90} \\
\midrule
SD-MCTS & 2442.2 & 100.4 & 23.1 \\
MPC-H4 & 2811.6 & 112.0 & 22.5 \\
MPC-H8 & 2588.6 & 108.8 & 16.0 \\
\bottomrule
\end{tabular}
\end{table}

\textbf{关键发现}：

\begin{enumerate}
\def\labelenumi{\arabic{enumi}.}
\item
  \textbf{最优MPC视界随环境变化}：稳态下H8最优，冲击下H4最优。这对MPC实践者提出了视界选择的两难困境。
\item
  \textbf{SD-MCTS跨环境鲁棒}：SD-MCTS在两种环境下均能达到接近最优MPC的性能，无需手动调优视界。
\item
  \textbf{尾部风险可控}：尽管MPC-H4均值成本与SD-MCTS相近，但其90分位成本（2811.6）显著高于SD-MCTS（2442.2），表明SD-MCTS在最差情况下更稳定。
\item
  \textbf{峰值积压优势明显}：相比MPC-H8，SD-MCTS将峰值积压降低25\%（大效应量），表明价值网络能更好地预测并缓冲需求冲击。
\end{enumerate}

\begin{quote}
【图占位】图5：需求冲击下的库存与积压动态轨迹 - 时间范围：\(t=0\) 至
\(t=60\) - 上图：各级库存水平（零售商、批发商、分销商、工厂），突出显示
\(t=15-30\) 冲击期 -
下图：各级积压水平，对比SD-MCTS（实线）与MPC-H8（虚线） -
关键标注：SD-MCTS峰值积压68.5 vs MPC-H8峰值91.5（降低25\%） -
阴影区域标记需求涌动期间（\(t=15-30\)，\(D=16\)）
\end{quote}

\begin{center}\rule{0.5\linewidth}{0.4pt}\end{center}

\textbf{实验3C：稳健性与管理情景压力测试 (Causal Intervention Stress
Tests)}

在运筹运营中，最具破坏性的风险并非参数在分布内的连续漂移，而是发生如``工厂意外停工（Capacity
Shock）''或``物流通信断链（Comm
Severed）''等极端底层物理结构突变。传统纯数据驱动的方法在这种
OOD（分布外）因果干预下面临极大概率发生崩溃，这是因为它们没有因果约束而仅仅学到了过拟合的相关性。然而，SD-MCTS
凭借可模块化编辑的守恒图结构先验，展现了卓越的抗突变韧性。

我们在测试推理阶段实施零样本（Zero-shot）的不可导人工干预： 1.
\textbf{产能冲击（工厂意外停工）}：在系统稳定期后，强行将底层工厂的最大产能从全量削减一半（即
\(C \to C/2\) 截断点强制干预）。 2.
\textbf{断链冲击（通信系统瘫痪）}：强行阻断分销商向下游节点传递的需求和订单数据长达
5 期。

\textbf{表2e：抗冲击压力测试结果 (相较于稳态最优解的性能贬损)}

\begin{table}[h]
\centering
\begin{tabular}{p{2.5cm}p{2.8cm}p{3.2cm}p{3cm}p{3cm}}
\toprule
\textbf{危机情境} & \textbf{SD-MCTS 成本升幅} & \textbf{SD-MCTS 峰值积压涨幅} & \textbf{MPC-H8 策略灾难崩溃率} & \textbf{端到端黑盒策略崩溃率} \\
\midrule
突发性停工 & +18.4\% & +33\% (受控缓冲) & 13\% & \textbf{100\% (极值积压至停摆)} \\
跨节点断链 & +27.6\% & +45\% (延迟传递) & 45\% (牛鞭爆震) & \textbf{100\% (零响应震荡)} \\
\bottomrule
\end{tabular}
\end{table}

\textbf{核心发现与政策启示}：
当底层运营环境的网络拓扑或规则物理突变时，依赖原分布强化学习重新训练的黑盒方案因缺乏图层结构修改手段将全面面临不可用硬伤；相较之下，拥有守恒规则先验网络逻辑层的
SD-MCTS
模型允许我们在规划源头上即插即用地只修改受损算子参数（如截断算子上限
\(C\)
被临时斩半）。它使得该框架能即时充当应对极端黑天鹅事件干预或开展情境应急演练工具（Crisis
Scenario Planning）高度可靠的``数字稳健决策沙盘''。

\begin{center}\rule{0.5\linewidth}{0.4pt}\end{center}

\hypertarget{ux673aux5236ux7814ux7a76}{%
\subsubsection{7.2 机制研究}\label{ux673aux5236ux7814ux7a76}}

我们进行受控实验以隔离各组件的贡献。

\hypertarget{ux6709ux6548ux89c6ux754cux6269ux5c55ux5b9aux74061ux9a8cux8bc1}{%
\paragraph{7.2.1
有效视界扩展（定理1验证）}\label{ux6709ux6548ux89c6ux754cux6269ux5c55ux5b9aux74061ux9a8cux8bc1}}

\textbf{实验设计}：使用\emph{相同}学习世界模型比较SD-MCTS和MPC(V=0)，仅变化终端值评估。

\textbf{表3：有效视界扩展 --- 场景2（10个种子）}

\begin{table}[h]
\centering
\begin{tabular}{p{4cm}p{3cm}p{3cm}p{3cm}}
\toprule
\textbf{规划视界 \(H\)} & \textbf{SD-MCTS} & \textbf{MPC (V=0)} & \textbf{视界差距} \\
\midrule
1 & \textbf{100\%} & 0\% & ≥7步 \\
2 & \textbf{100\%} & 0\% & ≥6步 \\
4 & \textbf{100\%} & 0\% & ≥4步 \\
6 & \textbf{100\%} & 0\% & ≥2步 \\
8 & \textbf{100\%} & 90\% & --- \\
10 & \textbf{100\%} & 100\% & 0 \\
\bottomrule
\end{tabular}
\end{table}

\textbf{关键发现}：\(H=1\) 的SD-MCTS达到MPC(V=0)需要 \(H \geq 8\)
才能匹配的效果------\textbf{有效视界扩展至少7步}。

这直接验证了定理1：价值网络将无限时域信息压缩为单个标量，使短深度搜索能够``看到''遥远的未来。

\hypertarget{ux6a21ux62dfux9884ux7b97ux654fux611fux6027}{%
\paragraph{7.2.2
模拟预算敏感性}\label{ux6a21ux62dfux9884ux7b97ux654fux611fux6027}}

\textbf{实验设计}：变化MCTS模拟次数 \(N_{\text{sim}}\)
并比较三种叶评估策略。

\textbf{表4：模拟预算敏感性 --- 场景2（30个种子）}

\begin{table}[h]
\centering
\begin{tabular}{p{4cm}p{3cm}p{3cm}p{3cm}}
\toprule
\textbf{模拟倍数 \(M\)†} & \textbf{SD-MCTS} & \textbf{MCTS-Zero} & \textbf{MCTS-Random} \\
\midrule
2 & \textbf{100\%} & 0\% & 0\% \\
6 & \textbf{100\%} & 6.7\% & 0\% \\
10 & \textbf{100\%} & 53.3\% & 0\% \\
14 & \textbf{100\%} & 90.0\% & 0\% \\
\bottomrule
\end{tabular}
\end{table}

†\(N_{\text{sim}} = \max(M \times 5, 20)\)

\textbf{关键发现}：SD-MCTS在所有预算下保持100\%成功率。MCTS-Zero（无价值网络）需要\textbf{7倍更多模拟}才能接近SD-MCTS性能。MCTS-Random（未训练网络）完全失败，确认性能源于\emph{训练}而非架构。

\textbf{启示}：投资于离线价值学习比投资于在线搜索预算产生更高回报。

\hypertarget{sd-mcts-vs.oracle-mpcux8ba1ux7b97ux6548ux7387}{%
\paragraph{7.2.3 SD-MCTS vs.~Oracle
MPC：计算效率}\label{sd-mcts-vs.oracle-mpcux8ba1ux7b97ux6548ux7387}}

\textbf{表5：计算比较 --- 场景2（10个种子）}

\begin{table}[h]
\centering
\begin{tabular}{p{2cm}p{3.5cm}p{3.5cm}p{3cm}p{3cm}}
\toprule
\textbf{视界 \(H\)} & \textbf{SD-MCTS成功率} & \textbf{Oracle MPC成功率} & \textbf{SD-MCTS时间} & \textbf{Oracle MPC时间} \\
\midrule
1 & \textbf{100\%} & 0\% & 57 ms & 344 ms \\
4 & \textbf{100\%} & 0\% & 112 ms & 669 ms \\
8 & \textbf{100\%} & 90\% & 184 ms & 1477 ms \\
10 & \textbf{100\%} & 100\% & 214 ms & 1812 ms \\
\bottomrule
\end{tabular}
\end{table}

\textbf{关键发现}：SD-MCTS以\textbf{Oracle MPC
3.2\%的计算成本}达到100\%成功率（等效性能下57 ms vs.~1812 ms）。

\begin{center}\rule{0.5\linewidth}{0.4pt}\end{center}

\hypertarget{ux53efux6269ux5c55ux6027ux548cux9c81ux68d2ux6027-ux5f85ux8865ux5145}{%
\subsubsection{7.3 可扩展性和鲁棒性
{[}待补充{]}}\label{ux53efux6269ux5c55ux6027ux548cux9c81ux68d2ux6027-ux5f85ux8865ux5145}}

以下实验计划用于回应审稿人关切：

\hypertarget{ux52a8ux4f5cux7a7aux95f4ux89c4ux6a21-ux8ba1ux5212ux4e2d}{%
\paragraph{7.3.1 动作空间规模
{[}计划中{]}}\label{ux52a8ux4f5cux7a7aux95f4ux89c4ux6a21-ux8ba1ux5212ux4e2d}}

\begin{table}[h]
\centering
\begin{tabular}{p{2cm}p{3cm}p{3.5cm}p{3.5cm}p{3cm}}
\toprule
\textbf{工人 \(N\)} & \textbf{动作空间 \(|A|\)} & \textbf{SD-MCTS回报} & \textbf{SD-MCTS时间} & \textbf{备注} \\
\midrule
2 & 3 & [现有] & [现有] & 基线 \\
5 & 6 & [待运行] & [待运行] & --- \\
10 & 11 & [待运行] & [待运行] & --- \\
20 & 21 & [待运行] & [待运行] & 可能需要渐进加宽 \\
\bottomrule
\end{tabular}
\end{table}

\hypertarget{ux53c2ux6570ux6f02ux79fb-oodux9c81ux68d2ux6027-ux8ba1ux5212ux4e2d}{%
\paragraph{7.3.2 参数漂移 / OOD鲁棒性
{[}计划中{]}}\label{ux53c2ux6570ux6f02ux79fb-oodux9c81ux68d2ux6027-ux8ba1ux5212ux4e2d}}

\begin{table}[h]
\centering
\begin{tabular}{p{3cm}p{3.5cm}p{3.5cm}p{4cm}}
\toprule
\textbf{漂移幅度} & \textbf{SD-MCTS} & \textbf{Oracle MPC} & \textbf{备注} \\
\midrule
0\%（分布内） & [现有] & [现有] & 基线 \\
±15\% & [待运行] & [待运行] & --- \\
±30\% & [待运行] & [待运行] & --- \\
\bottomrule
\end{tabular}
\end{table}

\hypertarget{ux7ea6ux675fux8fddux53cdux7edfux8ba1-ux8ba1ux5212ux4e2d}{%
\paragraph{7.3.3 约束违反统计
{[}计划中{]}}\label{ux7ea6ux675fux8fddux53cdux7edfux8ba1-ux8ba1ux5212ux4e2d}}

\begin{table}[h]
\centering
\begin{tabular}{p{5cm}p{3cm}p{3cm}p{3cm}}
\toprule
\textbf{方法} & \textbf{队列<0率} & \textbf{效率>1率} & \textbf{疲劳<0率} \\
\midrule
SD-MCTS（结构化） & [预期：0\%] & [预期：0\%] & [预期：0\%] \\
黑盒MCTS & [待运行] & [待运行] & [待运行] \\
\bottomrule
\end{tabular}
\end{table}

\begin{center}\rule{0.5\linewidth}{0.4pt}\end{center}

\hypertarget{ux5b9eux8df5ux6d1eux89c1}{%
\subsubsection{7.4 实践洞见}\label{ux5b9eux8df5ux6d1eux89c1}}

\hypertarget{ux89c6ux754c-ux7cbeux5ea6ux6743ux8861ux7684ux5b9eux8df5ux5e94ux7528}{%
\paragraph{7.4.1
视界-精度权衡的实践应用}\label{ux89c6ux754c-ux7cbeux5ea6ux6743ux8861ux7684ux5b9eux8df5ux5e94ux7528}}

我们将发现综合为实践者的可操作指南：

\textbf{定量误差分解}（来自定理2，§5.3）：

\begin{table}[h]
\centering
\begin{tabular}{p{4cm}p{6cm}p{4cm}}
\toprule
\textbf{误差来源} & \textbf{公式} & \textbf{典型值} \\
\midrule
MPC截断偏差 & \(\gamma^H R_{\max}/(1-\gamma)\) & \textasciitilde{}860 \\
SD-MCTS参数误差 & \(L_\theta \delta_\theta\) & \textasciitilde{}2.5 \\
SD-MCTS价值误差 & \(\gamma^H \epsilon_V\) & \textasciitilde{}4.3 \\
\midrule
\textbf{比率} & 截断 / SD-MCTS & \textbf{\textasciitilde{}126×} \\
\bottomrule
\end{tabular}
\end{table}

\textbf{启示}：在时间常数 \(\tau > H\)
的系统中，截断偏差以两个数量级占主导。投资于价值网络训练比投资于参数校准产生远更高的回报。

\hypertarget{ux76f8ux56fesd-mctsux4f55ux65f6ux5360ux4f18ux5f85ux8865ux5145}{%
\paragraph{7.4.2
相图：SD-MCTS何时占优？{[}待补充{]}}\label{ux76f8ux56fesd-mctsux4f55ux65f6ux5360ux4f18ux5f85ux8865ux5145}}

\begin{quote}
【图占位】图6：相图 - X轴：\(H/\tau\)（规划视界 / 系统时间常数） -
Y轴：\(\delta_\theta\)（参数误差） - 区域：SD-MCTS优势（蓝色）
vs.~Oracle MPC优势（红色） - 叠加：实验的经验数据点
\end{quote}

\textbf{初步洞见}：SD-MCTS在
\(H/\tau < 1\)（规划视界短于系统时间常数）时占优。这描述了大多数具有延迟效应的OR系统：疲劳、磨损、季节性需求、库存累积。

\hypertarget{ux5173ux952eux7ed3ux679cux603bux7ed3}{%
\paragraph{7.4.3
关键结果总结}\label{ux5173ux952eux7ed3ux679cux603bux7ed3}}

\begin{table}[h]
\centering
\begin{tabular}{p{5cm}p{5cm}p{5cm}}
\toprule
\textbf{发现} & \textbf{证据} & \textbf{实践启示} \\
\midrule
学习参数 + 深度规划 > 完美参数 + 浅层规划 & 表1：比Oracle MPC改进3.0\% & 投资价值学习，而非参数校准 \\
价值网络扩展有效视界≥7步 & 表3：H=1匹配H=8 & 有好的终端值时短在线搜索足够 \\
价值学习减少搜索预算需求7倍 & 表4：SD-MCTS对模拟次数鲁棒 & 离线训练摊销在线计算 \\
SD-MCTS以3.2\%计算成本达到相同性能 & 表5：57 ms vs. 1812 ms & 实时部署可行 \\
供应链场景下性能媲美最优MPC视界 & 表2c：SD-MCTS (623.6) ≈ MPC-H8 (580.9) & 无需手动调优规划视界 \\
\textbf{最优MPC视界随环境动态变化} & 表2d：稳态H8最优，冲击下H4最优 & MPC需针对不同场景调优视界 \\
SD-MCTS跨环境鲁棒，无需视界调优 & 表2d：稳态≈H8，冲击≈H4 & 单一模型适应多种运营条件 \\
需求冲击下峰值积压降低25\%（vs H8） & 表2d：68.5 vs 91.5 & 提升供应链韧性的实用价值 \\
\bottomrule
\end{tabular}
\end{table}

\begin{center}\rule{0.5\linewidth}{0.4pt}\end{center}

\hypertarget{ux53efux6269ux5c55ux6027ux8f6fux4ef6ux548cux6570ux636e}{%
\subsection{8.
可复现性、软件和数据}\label{ux53efux6269ux5c55ux6027ux8f6fux4ef6ux548cux6570ux636e}}

根据IJOC可复现性政策，我们提供完整的软件包以实现所有报告结果的独立验证。

\hypertarget{ux4ed3ux5e93ux7ed3ux6784}{%
\subsubsection{8.1 仓库结构}\label{ux4ed3ux5e93ux7ed3ux6784}}

代码仓库遵循IJOC推荐的组织方式：

\begin{verbatim}
sd-mcts-ijoc/
|-- README.md              # 快速入门指南和复现说明
|-- LICENSE                # MIT许可证
|-- AUTHORS                # 作者信息和贡献
|-- requirements.txt       # Python依赖及版本
|
|-- src/                   # 源代码
|   |-- env/               # 环境实现
|   |   |-- task_allocation_env.py
|   |   |-- delayed_collapse_env.py
|   |   +-- beer_game_env.py
|   |-- world_model/       # 结构化世界模型
|   |   |-- task_sd_model.py       # 任务分配SD模型
|   |   |-- delayed_collapse_model.py  # 延迟崩溃SD模型
|   |   +-- beer_game_sd_model.py  # 啤酒游戏SD模型
|   |-- planning/          # MCTS实现
|   |   +-- mcts.py
|   |-- agent/             # 神经网络
|   |   +-- networks.py    # 策略网络和价值网络
|   +-- baselines/         # 领域基线策略
|       +-- beer_game_baselines.py
|
|-- experiments/           # 实验脚本
|   |-- baselines/         # MPC基线实现
|   |   |-- oracle_mpc.py
|   |   |-- adaptive_mpc.py
|   |   +-- rls_estimator.py
|   |-- run_m1.py          # 场景1-2实验
|   +-- run_beer_game_v2.py # 场景3：啤酒游戏实验
|
|-- results/               # 原始实验输出
|   |-- scenario1/         # 场景1日志和指标
|   |-- scenario2/         # 场景2日志和指标
|   +-- beer_game_v2/      # 啤酒游戏主实验
|       |-- horizon_wall/  # 实验3A结果
|       +-- resilience/    # 实验3B结果
|
|-- ijoc_paper_draft.md    # 英文论文
+-- ijoc_paper_draft_zh.md # 中文论文
\end{verbatim}

\hypertarget{ux590dux73b0ux8bf4ux660e}{%
\subsubsection{8.2 复现说明}\label{ux590dux73b0ux8bf4ux660e}}

\hypertarget{ux5febux901fux5f00ux59cb}{%
\paragraph{快速开始}\label{ux5febux901fux5f00ux59cb}}

\begin{verbatim}
# 克隆仓库
git clone https://github.com/[anonymized]/sd-mcts-ijoc.git
cd sd-mcts-ijoc

# 创建环境
conda create -n sdmcts python=3.10
conda activate sdmcts
pip install -r requirements.txt

# 复现主要结果
python experiments/run_m1.py --seeds 50           # 场景1-2
python experiments/run_beer_game_v2.py --horizon-wall --eval-episodes 30  # 实验3A
python experiments/run_beer_game_v2.py --resilience --env-steps 10000 --eval-episodes 30  # 实验3B
\end{verbatim}

\hypertarget{ux590dux73b0ux7279ux5b9aux8868ux683c}{%
\paragraph{复现特定表格}\label{ux590dux73b0ux7279ux5b9aux8868ux683c}}

\begin{table}[h]
\centering
\begin{tabular}{p{3cm}p{6cm}p{6cm}}
\toprule
\textbf{表格} & \textbf{命令} & \textbf{输出文件} \\
\midrule
表1 (场景1) & \texttt{python experiments/run\_m1.py --seeds 50} & \texttt{results/scenario1/} \\
表2 (场景2) & \texttt{python experiments/run\_m1.py --scenario 2 --seeds 10} & \texttt{results/scenario2/} \\
表2c (视界墙) & \texttt{python experiments/run\_beer\_game\_v2.py --horizon-wall --eval-episodes 30} & \texttt{results/beer\_game\_v2/horizon\_wall/horizon\_wall\_results.json} \\
表2d (韧性) & \texttt{python experiments/run\_beer\_game\_v2.py --resilience --env-steps 10000 --eval-episodes 30} & \texttt{results/beer\_game\_v2/resilience/resilience\_results.json} \\
表3 (视界扩展) & \texttt{python experiments/run\_m1.py --horizon-ablation} & \texttt{results/horizon\_ablation/} \\
\bottomrule
\end{tabular}
\end{table}

\hypertarget{ux786cux4ef6ux548cux8f6fux4ef6ux73afux5883}{%
\subsubsection{8.3
硬件和软件环境}\label{ux786cux4ef6ux548cux8f6fux4ef6ux73afux5883}}

\textbf{硬件}（用于报告结果）： - CPU：Intel Xeon E5-2680 v4 @ 2.40GHz -
内存：64 GB - GPU：不需要（仅CPU）

\textbf{软件}： - Python 3.9.7 - PyTorch 1.12.0 - NumPy 1.21.0 -
完整列表见 \texttt{requirements.txt}

\hypertarget{ux6570ux636eux53efux7528ux6027}{%
\subsubsection{8.4 数据可用性}\label{ux6570ux636eux53efux7528ux6027}}

本研究仅使用\textbf{合成环境}；不需要外部数据集。环境动态在§6.2中完整指定，并在
\texttt{src/env/} 中实现。

所有实验的随机种子已固定并记录，确保精确可复现性。

\begin{center}\rule{0.5\linewidth}{0.4pt}\end{center}

\hypertarget{ux7ed3ux8bba}{%
\subsection{9. 结论}\label{ux7ed3ux8bba}}

我们介绍了SD-MCTS，一个神经-符号规划框架，将来自系统动力学的结构化世界模型与用于无限时域近似的学习价值函数相结合。本文的核心设计原则可以概括为：

\begin{quote}
\textbf{``结构实现样本效率，价值实现视界扩展''（Structure for Sample
Efficiency, Value for Horizon Extension）}
\end{quote}

结构化世界模型约束假设空间，解决OR典型的``小数据''问题（\(N \approx 10^4\)）；学习价值网络压缩无限时域尾部，解决MPC的``短视界''问题（\(\tau > H\)）。我们的工作表明，基于模型和无模型方法之间的传统二分法存在一个富有成效的中间地带：\textbf{结构化学习}，其中领域知识约束模型形式，而数据决定模型参数。我们期望在能源、物流和医疗等对规划误差或幻觉（Hallucination）零容忍的安全敏感领域，SD-MCTS 能够提供一个完全可审计且\textbf{零幻觉的坚实物理底座}。综上所述，我们将此框架视作迈向\textbf{广义 Neuro-OR} 范式的重要一步，其核心在于让数据驱动代理在一个由\textbf{物理定律设定的导轨}上安全且极简地行驶，而非落入完全黑盒状态下的发散搜索。

\hypertarget{ux6838ux5fc3ux6d1eux89c1}{%
\subsubsection{9.1 核心洞见}\label{ux6838ux5fc3ux6d1eux89c1}}

四个主要发现从本工作中涌现，我们按重要性排序呈现：

\textbf{洞见1（反传统发现）：在延迟效应系统中，规划深度主导参数精度。}
我们最具颠覆性的结果：\emph{学习}参数的SD-MCTS优于\emph{完美}参数的Oracle
MPC（表1：−7853 vs −8098，\(p < 0.0001\)，Cohen's
\(d = 2.05\)）。这挑战了传统OR智慧``先建好模型再控制''的假设。\textbf{边界条件}：当系统时间常数超过规划视界（\(\tau > H\)）时，视界截断误差以两个数量级主导参数不精确误差（§7.4.1）。这为实践者提供了明确的\textbf{决策边界}：当延迟效应的时间尺度超过MPC规划视界2倍以上时，应优先投资于价值函数近似而非穷尽参数校准。

\textbf{洞见2：结构化模型实现样本效率。} 在OR典型的 \(\sim 10^4\)
交互约束下，结构化世界模型实现了无模型方法无法做到的：具有可量化不确定性的可靠策略优化。存量-流量形式提供归纳偏置，将样本复杂度从
\(O(|\mathcal{S}||\mathcal{A}|/\epsilon^2)\) 降低到
\(O(d/\epsilon^2)\)，其中 \(d\)
是可学习参数维度（引理1，在表1-2中验证）。

\textbf{洞见3：价值网络实现有效视界扩展。} 学习终端成本 \(V_\psi\)
将有效规划视界扩展7+步（表3：\(H=1\) 的SD-MCTS匹配 \(H \geq 8\)
的MPC），使短深度MCTS能够以3.2\%的计算成本匹配长视界MPC的性能。这验证了ADP原则：近似价值函数可以替代显式多步前瞻。

\textbf{洞见4（洞见3的因果结果）：SD-MCTS对视界选择具有鲁棒性。}
由于价值网络将性能与显式搜索深度\textbf{解耦}，SD-MCTS在不同规划视界下表现稳定（\(H \in [1, 14]\)
范围内性能退化
\(< 1\%\)）。在啤酒游戏实验中，MPC的最优视界随环境条件变化（稳态H8最优，冲击下H4最优），而SD-MCTS无需重新调优即可在两种条件下达到接近最优性能。此外，相比固定视界的MPC-H8，SD-MCTS在需求冲击下将峰值积压降低25\%（Cohen's
\(d = 0.90\)），表明价值网络能够预测并缓冲供应链扰动。

\textbf{洞见5（管理情境韧性）：可编辑框架能充当应对极端破坏性危机的沙盘演练工具。}
在诸如突发供需切断、基础物理设施半毁等脱离正常训练分布域的外界不可抗干预中（OOD
因果破坏），脱离了严谨拓扑架构的数据驱动黑盒的机制呈现了 100\%
策略灾难崩溃。相反地，拥有硬编码的可手动编辑物理包络参数的 SD-MCTS
框架能在遭遇真实世界物理结构突变的第一时间，以零等待成本（无需收集受损后的轨迹重训网络，反而仅需根据灾损参数就地修改守恒限制方程组常量）迅猛切换出合规的避险转移路线。这体现出混合型强化学习具备传统黑盒模型难以企及的高度管理学韧性，它是迈向更优越的安全运筹学智能控制的关键一步。

\textbf{洞见6（实时计算可行性）：极大削减的在线计算开销（Computational Overhead）。}
除了在数据效率上的突破，物理先验极大缩减了参数空间带来的另一重优势，凸显在在线控制周期（Real-time Control Cycle）内的计算经济性。面对真实的生产线与物流频控，算法必须在边缘级算力下低延迟运作。主流世界模型（如拥军千万级权重的 DreamerV3/RSSM）在推理时需要庞大的矩阵相乘吞吐和梯度反传算力。而 SD-MCTS 的核心差分世界模型仅囊括数十个一维物理参量（\(\sim \mathcal{O}(10)\)），更因剥离了无结构的深层盲目搜索，将截断视界压至极浅（\(H=1\)）从而避免了自发扩展搜索树节点的指数级开销。这种在线微推断与网络轻量双管齐下的轻算特质，使其彻底具备了替代传统启发式在嵌入式强实时系统中执行计算的可行基础。

\hypertarget{ux5b9eux8df5ux5efaux8bae}{%
\subsubsection{9.2 实践建议}\label{ux5b9eux8df5ux5efaux8bae}}

对于考虑SD-MCTS或相关方法的实践者，我们提供以下指导：

\textbf{何时使用SD-MCTS：} -
具有已知物理/运营结构的系统（守恒律、流量约束） - 时间常数 \(\tau > H\)
的延迟反馈效应 - 样本预算 \(< 10^5\)（无模型方法困难的地方） -
对可解释、可审计决策的需求

\textbf{何时偏好替代方案：} - 动态未知的黑盒系统 -
数据充足（\textgreater{}\(10^6\) 样本）时无模型方法表现出色 -
非常短的时间常数（\(\tau \ll H\)）时MPC足够

\textbf{资源配置指导：}
给定固定计算预算，投资于价值网络训练而非参数校准。25\%参数误差的训练良好价值网络优于无终端值近似的完美参数。

\hypertarget{ux5c40ux9650ux6027}{%
\subsubsection{9.3 局限性}\label{ux5c40ux9650ux6027}}

我们承认几个局限性：

\begin{enumerate}
\def\labelenumi{\arabic{enumi}.}
\item
  \textbf{结构-参数分离假设。}
  我们的框架假定了一个基本正确的存量-流量因果骨架。在实践中，如果结构部分未知或存在先验偏误（Misspecification），可通过系统性观测残差来察觉。未来可考虑整合可微因果发现（Differentiable Causal Discovery，如 NOTEARS）或时序系统辨识算法（如 PCMCI），实现由数据驱动的物理图谱自动纠偏与边缘关系推断。
\item
  \textbf{连续松弛。}
  用于可微分学习的平均场近似可能在高度离散系统中引入偏差。
\item
  \textbf{规模。} 当前实验使用 \(N \leq 2\) 工人；扩展到 \(N > 10\)
  需要渐进加宽或动作抽象。
\item
  \textbf{领域特异性。}
  虽然我们的理论框架是通用的，但实证验证聚焦于劳动力管理场景。
\item
  \textbf{参数不确定性与模型失配。}
  当前框架假设参数固定（尽管未知），且在规划期间不对参数不确定性或漂移建模。实践中，参数可能随时间变化（如到达率季节性波动、疲劳敏感度随劳动力组成变化、设备退化随时间加速）。处理非平稳性需要显式机制，例如：

  \begin{itemize}
  \tightlist
  \item
    \textbf{集成规划}：维护参数估计的分布，在多个场景上规划
  \item
    \textbf{在线参数适应}：使用贝叶斯更新或递归最小二乘跟踪漂移参数
  \item
    \textbf{场景树MCTS}：在树搜索过程中对参数实现进行分支（Couëtoux等人
    2011）
  \end{itemize}

  这些扩展超出当前范围，留待未来工作。
\end{enumerate}

\hypertarget{ux672aux6765ux5de5ux4f5c}{%
\subsubsection{9.4 未来工作}\label{ux672aux6765ux5de5ux4f5c}}

本工作开启了几个研究方向：

\textbf{方法论扩展：} -
\textbf{自动结构发现}：从数据中学习存量-流量拓扑，可能通过因果发现方法 -
\textbf{层次规划}：将大系统分解为具有学习接口的交互子系统 -
\textbf{鲁棒优化}：扩展到参数不确定性被显式建模的分布鲁棒设置

\textbf{应用领域：} -
\textbf{排队网络}：存量-流量表示的自然契合，具有经过验证的运筹学价值 -
\textbf{Petri网}：具有令牌守恒律的离散事件系统 -
\textbf{供应链动态}：具有牛鞭效应的多层库存 -
\textbf{流行病建模}：具有干预规划的分区模型（SIR/SEIR）

\textbf{走向Neuro-OR：}
我们设想一个更广泛的研究项目------``Neuro-OR''------结合运筹学模型的结构保证与神经网络的函数近似能力。SD-MCTS代表这个设计空间中的一个点；许多其他点仍未探索。关键洞见是\textbf{领域知识和数据是互补的，而非替代的}：结构提供样本效率和可解释性；学习提供适应性和近似能力。

\begin{center}\rule{0.5\linewidth}{0.4pt}\end{center}

\hypertarget{ux9644ux5f55aux7b26ux53f7ux6c47ux603b}{%
\subsection{附录A：符号汇总}\label{ux9644ux5f55aux7b26ux53f7ux6c47ux603b}}

\begin{table}[h]
\centering
\begin{tabular}{p{3cm}p{11cm}}
\toprule
\textbf{符号} & \textbf{定义} \\
\midrule
\(s, s_t\) & 状态（通用 / 时刻 \(t\)） \\
\(a, a_t\) & 动作（通用 / 时刻 \(t\)） \\
\(\theta, \hat{\theta}\) & 真实 / 估计世界模型参数 \\
\(V^*, V_\psi\) & 最优 / 学习价值函数 \\
\(H\) & 规划视界（MCTS深度） \\
\(\tau\) & 系统时间常数（示例中的疲劳阈值） \\
\(N_{\text{sim}}\) & MCTS模拟次数 \\
\(\gamma\) & 折扣因子 \\
\bottomrule
\end{tabular}
\end{table}

\begin{center}\rule{0.5\linewidth}{0.4pt}\end{center}

\hypertarget{ux9644ux5f55bux8bc1ux660e}{%
\subsection{附录B：证明}\label{ux9644ux5f55bux8bc1ux660e}}

\hypertarget{b.1-ux5f15ux74061ux7684ux8bc1ux660eux6837ux672cux590dux6742ux5ea6}{%
\subsubsection{B.1
引理1的证明（样本复杂度）}\label{b.1-ux5f15ux74061ux7684ux8bc1ux660eux6837ux672cux590dux6742ux5ea6}}

\textbf{线性参数化假设的精确表述}：

设 \(f: \mathcal{S} \times \mathcal{A} \times \Theta \to \mathcal{S}\)
为世界模型的确定性部分。核心假设为：

\begin{enumerate}
\def\labelenumi{\arabic{enumi}.}
\tightlist
\item
  \textbf{线性参数化}：\(f(s,a;\theta) = \Phi(s,a)\theta\)，其中
  \(\Phi: \mathcal{S} \times \mathcal{A} \to \mathbb{R}^{n \times d}\)
  为已知特征映射
\item
  \textbf{特征有界}：\(\|\Phi(s,a)\|_2 \leq B\) 对所有 \((s,a)\) 成立
\item
  \textbf{持续激励}：\(\lambda_{\min}(\mathbb{E}[\Phi^\top \Phi]) \geq \lambda_0 > 0\)
\item
  \textbf{亚高斯噪声}：\(\|\xi\|_{\psi_2} \leq \sigma\)
\end{enumerate}

\textbf{完整证明（基于集中不等式）}：

\emph{Step 1}：最小二乘估计的正规方程：

\[
\hat{\theta}_N = \left(\sum_{i=1}^N \Phi_i^\top \Phi_i\right)^{-1} \sum_{i=1}^N \Phi_i^\top s'_i
\]

其中 \(s'_i = \Phi_i \theta^* + \xi_i\)。代入得：

\[
\hat{\theta}_N - \theta^* = \left(\sum_{i=1}^N \Phi_i^\top \Phi_i\right)^{-1} \sum_{i=1}^N \Phi_i^\top \xi_i
\]

\emph{Step 2}：由持续激励和 Hoeffding 不等式，以高概率：

\[
\frac{1}{N}\sum_{i=1}^N \Phi_i^\top \Phi_i \succeq \frac{\lambda_0}{2} I_d
\]

\emph{Step 3}：对 \(\sum_i \Phi_i^\top \xi_i\) 应用向量 Bernstein
不等式。由于 \(\xi_i\) 亚高斯且与 \(\Phi_i\) 独立：

\[
\mathbb{P}\left(\left\|\sum_{i=1}^N \Phi_i^\top \xi_i\right\|_2 > t\right) \leq 2d \exp\left(-\frac{t^2}{2\sigma^2 B^2 N}\right)
\]

取 \(t = \sigma B \sqrt{2dN\log(2d/\delta)}\)，以概率至少 \(1-\delta\)：

\[
\left\|\sum_i \Phi_i^\top \xi_i\right\|_2 \leq O(\sigma B\sqrt{dN\log(d/\delta)})
\]

\emph{Step 4}：综合 Step 2 和 Step 3：

\[
\|\hat{\theta}_N - \theta^*\|_2 \leq \frac{2}{N\lambda_0} \cdot O(\sigma B\sqrt{dN\log(d/\delta)}) = O\left(\sqrt{\frac{d\log(d/\delta)}{N\lambda_0}}\right)
\]

令此式 \(\leq \epsilon\)，解得
\(N = O(d\log(d/\delta)/(\lambda_0\epsilon^2))\)。

\hypertarget{b.2-ux5b9aux74061ux7684ux8bc1ux660eux6709ux6548ux89c6ux754cux6269ux5c55}{%
\subsubsection{B.2
定理1的证明（有效视界扩展）}\label{b.2-ux5b9aux74061ux7684ux8bc1ux660eux6709ux6548ux89c6ux754cux6269ux5c55}}

\textbf{模型误差累积的精确界}：

设世界模型单步误差为
\(\|f(s,a;\hat{\theta}) - f(s,a;\theta^*)\| \leq \epsilon_{\text{model}}\)。

经过 \(H\) 步展开，状态预测误差满足：

\[
\|\hat{s}_H - s_H\| \leq \sum_{t=0}^{H-1} L^{H-1-t} \epsilon_{\text{model}} = \epsilon_{\text{model}} \cdot \frac{L^H - 1}{L - 1}
\]

其中 \(L\) 是转移函数的 Lipschitz 常数。当
\(L \leq 1\)（收缩映射）时，此界简化为
\(H \cdot \epsilon_{\text{model}}\)。

\textbf{价值误差的完整分解}：

\[
|V^*(s) - \hat{V}_{\text{MCTS}}(s)| \leq \sum_{t=0}^{H-1} \gamma^t |r(s_t,a_t) - r(\hat{s}_t,\hat{a}_t)| + \gamma^H |V^*(s_H) - V_\psi(\hat{s}_H)|
\]

设奖励函数 \(L_r\)-Lipschitz，价值函数 \(L_V\)-Lipschitz：

\[
\leq L_r \sum_{t=0}^{H-1} \gamma^t \|\hat{s}_t - s_t\| + \gamma^H (L_V \|\hat{s}_H - s_H\| + \epsilon_V)
\]

代入状态误差累积界，得最终结论。

\hypertarget{b.3-ux5b9aux74062ux7684ux8bc1ux660eux89c6ux754c-ux7cbeux5ea6ux6743ux8861}{%
\subsubsection{B.3
定理2的证明（视界-精度权衡）}\label{b.3-ux5b9aux74062ux7684ux8bc1ux660eux89c6ux754c-ux7cbeux5ea6ux6743ux8861}}

\textbf{参数敏感性的推导}：

考虑 Bellman
方程：\(V^*(s;\theta) = \max_a \{R(s,a) + \gamma \mathbb{E}_{P_\theta}[V^*(s';\theta)]\}\)

对 \(\theta\) 求导（包络定理）：

\[
\nabla_\theta V^*(s;\theta) = \gamma \mathbb{E}_{P_\theta}[\nabla_\theta \log P_\theta(s'|s,a^*) \cdot V^*(s';\theta)] + \gamma \mathbb{E}_{P_\theta}[\nabla_\theta V^*(s';\theta)]
\]

取范数并递归展开 \(H\) 步：

\[
\|\nabla_\theta V^*(s;\theta)\| \leq \sum_{t=0}^{H-1} \gamma^t \|\nabla_\theta \log P_\theta\| \cdot \|V^*\|_\infty + \gamma^H \|\nabla_\theta V^*\|
\]

设
\(\|\nabla_\theta \log P_\theta\| \leq G\)，\(\|V^*\|_\infty \leq R_{\max}/(1-\gamma)\)，则：

\[
\|\nabla_\theta V^*\| \leq \frac{G \cdot R_{\max}}{(1-\gamma)^2}
\]

由 Lipschitz
条件：\(|V^*(s;\theta^*) - V^*(s;\hat{\theta})| \leq L_\theta \|\theta^* - \hat{\theta}\|\)，其中
\(L_\theta = \frac{G \cdot R_{\max}}{(1-\gamma)^2}\)。

\begin{center}\rule{0.5\linewidth}{0.4pt}\end{center}

\end{document}
